\documentclass[reprint,amsmath,amssymb,aps,prd,longbibliography,floatfix]{revtex4-2}
\usepackage{graphicx}
\usepackage{dcolumn}
\usepackage{bm}
\usepackage{hyperref}
\usepackage{booktabs}
\usepackage{siunitx}
\usepackage{multirow}
\usepackage[dvipsnames]{xcolor}
\usepackage{longtable}
% Force longtable captions to span the full text width (both columns),
% otherwise revtex's reprint mode renders them at column width.
\setlength{\LTcapwidth}{\textwidth}
\sisetup{
    range-units = single,
    separate-uncertainty = true,
    table-align-uncertainty = true,
}

\begin{document}

\title{Substrate Framework: Falsifiable Observational Predictions}

\author{Eliot Gable}
\email{10015996+egable@users.noreply.github.com}
\affiliation{Independent Researcher}

\date{2026-05-05}

\begin{abstract}
\noindent
A live audit of the substrate-framework's quantitative predictions
across particle, gravitational, and cosmological observables. Each
predicted value is paired with the current best measurement and the
deviation, anchored on three independent reading paths (CODATA
$m_e$, CODATA $m_\mu$~\cite{CODATA2022}, and the framework's own
Higgs vev $\approx 246.21965$~GeV). All values computed at
$\mathrm{mp.prec}=500$. Status flags: $\dagger$ tension, $\ddagger$
open prediction, $\dagger\dagger$ regime-tension. As of 2026-05-05,
no entry is in falsifying tension ($\geq 3\sigma$) with observation.
Three entries ($\rho_\Lambda$, $H_0$ tension magnitude, TOV bound)
are flagged~$\ddagger$ pending closure of upstream prediction work;
38 distinct observables are auto-flagged blue (estimate-pending);
14 half-integer-slot dark-mode candidates are auto-flagged violet
(mass predicted, but the framework cannot yet decide whether the
slot is occupied). Cell counts in Tables~1--2 exceed observable
counts because each observable is rendered across three anchor
columns. This document is intended as a falsification target.
\end{abstract}

%% ============================================================================
%% PLAIN-TEXT ABSTRACT (LaTeX-macro-free; for arXiv submission portal etc.)
%% ============================================================================
\iffalse

A live audit of the substrate-framework's quantitative predictions across
particle, gravitational, and cosmological observables. Each predicted value
is paired with the current best measurement and the deviation, anchored on
three independent reading paths (CODATA m_e, CODATA m_mu, framework Higgs
vev approx 246.21965 GeV). All values computed at mp.prec = 500. Status
flags: dagger tension, double-dagger open prediction, dagger-dagger regime-tension.
As of 2026-05-05, no entry is in falsifying tension with observation;
three entries (rho_Lambda, H_0 tension magnitude, TOV bound) are flagged
double-dagger pending closure of upstream prediction work; 38 observables
are auto-flagged blue (estimate-pending); 14 half-integer-slot dark-mode
candidates are auto-flagged violet (mass predicted but slot occupancy
unresolved). This document is intended as a falsification target.

\fi

\maketitle

\section{Methodology}

Predictions are computed from a single mass-anchor parameter
$h_{\rm eff}$ that is independently fixed three ways: from CODATA
$m_e$~\cite{CODATA2022}, from CODATA $m_\mu$, and from the
framework's own Higgs vev. The three anchor reading paths agree
to $\sim 158$~ppt (max pairwise relative deviation across the full
42-particle catalogue), with the spread carried mostly by the
e-anchor-vs-vev-anchor pair ($\sim 158$~ppt) and the
$\mu$-anchor-vs-vev-anchor pair ($\sim 153$~ppt); the
e-anchor-vs-$\mu$-anchor pair agrees at $\sim 4$~ppt. All
computations are performed at $\mathrm{mp.prec}=500$ in mpmath; no
floating-point compounding occurs.

\paragraph{What each anchor column means.}
The three anchor columns in Tables~1 and~2 are not three
independent measurements; each carries a specific epistemic status:
\begin{itemize}\itemsep0pt
\item The \textbf{electron-anchor} column inherits the
electron-mass CODATA central value as its base reading; every
prediction in that column is propagated from that single
experimental input. The column's residual against any other
observable is therefore bounded below by the CODATA $m_e$
experimental uncertainty (and above by the framework's structural
precision).
\item The \textbf{muon-anchor} column does the same with the
muon-mass CODATA central value as its single experimental input.
\item The \textbf{vev-anchor} column is qualitatively different:
it carries the framework's geometry-intrinsic prediction of the
Higgs vev as its base reading, with no CODATA particle-mass input.
The values in that column are pure framework predictions; their
residual against observation is not bounded below by any
experimental uncertainty band on a measured anchor.
\end{itemize}
A reader comparing the three columns is therefore comparing two
experimentally-anchored readings against one geometry-only
reading; agreement of all three to $\sim 4$~ppt is consistency
between two independent CODATA inputs and one parameter-free
geometric prediction.

\paragraph{Auto-flag legend (appears on every cell).} Each predicted
value's deviation against the observed value is automatically
compared against the observation's published experimental
uncertainty $\sigma_{\rm obs}$ (covering $\sim 70$ observables):
\begin{itemize}\itemsep0pt
\item \textcolor{OliveGreen}{$\checkmark$ \textbf{green:}}
$|{\rm dev}| \leq 1\,\sigma_{\rm obs}$ --- prediction sits within
the experimental band.
\item \textcolor{orange}{$\bullet$ \textbf{orange:}}
$1\,\sigma_{\rm obs} < |{\rm dev}| \leq 3\,\sigma_{\rm obs}$ ---
marginal.
\item \textcolor{red}{$\times$ \textbf{red:}}
$|{\rm dev}| > 3\,\sigma_{\rm obs}$ --- outside the experimental
band; falsifying tension if reproducible.
\item \textcolor{blue}{$?$ \textbf{blue:}} estimate-pending; the
prediction depends on upstream work that is not yet closed, so
the deviation is \emph{not yet falsifiable}. The work-remaining
tally (\S~Work remaining) names which item gates each blue
prediction.
\item \textcolor{violet}{$?$ \textbf{violet:}} mass predicted, but
the framework cannot yet decide whether this candidate slot is
physically occupied. Used for spin-$\tfrac12$ half-integer-slot
dark-mode candidates.
\item plain (no colour): $\sigma_{\rm obs}$ unavailable for this
observable; status undetermined.
\end{itemize}

\section{Cross-anchor consistency}

\input{tables/cross_anchor_summary}%

The framework's mass scale $h_{\rm eff}$ can be fixed from any one
of several \emph{dimensionally distinct} observables. The seven
anchors used in Table~\ref{tab:cross-anchor} and the companion data
file are not seven variations on the same input; they span a range
of physical dimensions:
\begin{itemize}\itemsep0pt
\item $m_e$ --- mass, $[M]$ (CODATA $m_e c^2$ in MeV);
\item $v_C$ --- mass, $[M]$ (G$_F$-derived Higgs vev in GeV);
\item $m_P c^2$ --- mass, $[M]$ at the Planck scale (CODATA, in GeV);
\item $G$ --- $[L^3 M^{-1} T^{-2}]$ (CODATA Newton's constant in
SI units);
\item $\lambda_{C,\tau}$ --- length, $[L]$ (PDG-derived Compton
wavelength of the tau in fm);
\item $\alpha_{\rm EM}$ --- dimensionless (CODATA fine-structure
constant);
\item $v_C^{\rm struct}$ --- mass, $[M]$ (the framework's own Higgs
vev, computed from $E_8$ invariants alone with \emph{no} PDG input;
included as a 7th column in the data file).
\end{itemize}
A consistency claim across these inputs is therefore not a
normalisation property: an internally-rigid framework reproduces a
mass-dimension prediction equally well from a length-dimension
input, a coupling-strength input, and a gravitational input.

\paragraph{Two distinct empirical claims.}
The cross-anchor exercise establishes \emph{rigidity} (anchor
independence) and is logically separate from \emph{accuracy} (the
match against observation, addressed by the prediction tables in
Sections~III--IV). Both must hold for the framework to be
correct; here we report the rigidity result, and rely on the
prediction tables for the accuracy result.

\paragraph{Rigidity result, operational.}
For every one of the \crossanchorNObservables{} predicted
observables in the framework's structural catalogue, the seven
anchor-derived predictions agree to \emph{better than
\crossanchorMaxPpm~ppm} (max pairwise spread), with a median
pairwise spread of \crossanchorMedianPpm~ppm. This consistency is a
structural property of the framework's equations: no anchor is
privileged, no observed value is tuned, and the residual spread
reflects the precision of the underlying CODATA/PDG inputs (PDG
$m_\tau$ at $\sim 6.8\times 10^{-5}$ and CODATA $G$ at $\sim
2.2\times 10^{-5}$ dominate) rather than tunable parameters. The
parameter-free $v_C^{\rm struct}$ anchor reproduces the
mass-anchored predictions at $<\!1$~ppb spread against $v_C^{\rm
obs}$, so the rigidity result holds even when the seventh
column carries no measured input at all.

\paragraph{Why this isn't dimensional analysis through known physics.}
A reader might ask: ``if every anchor produces the same prediction,
isn't that just because the anchors are related through known
physics?'' The framework's structural equations do \emph{not} route
through CODATA conversions. Each anchor enters its own
geometry-intrinsic relation: $m_e$ via the gateway-slot $R$-ratio
($m_e = R_e \cdot h_{\rm eff}$ with $R_e = (\dim - 2)/\dim$), $v_C$
via the chain-position ratio $R_v$, $m_P$ via the structural
factor $F_{\rm struct} = m_e/m_P$ derived from $E_8$ invariants
alone, $G$ via $F_{\rm struct}^2$ entering as a closed
geometric identity, $\lambda_{C,\tau}$ via the tau slot's
sub-harmonic position $R_\tau$, and $\alpha_{\rm EM}$ via the
lepton-vacuum-polarisation-corrected $\alpha_2 \sin^2\theta_W$
combination. The fact that all seven paths converge on the same
predicted value is a property of the framework's structural
identities (companion paper for the full derivation), not of
unit conversions; SI dimensional analysis cannot, on its own,
relate $G$ to $m_e$ to a part-per-million.

\begin{widetext}
\input{tables/table_cross_anchor}
\end{widetext}

\noindent\textbf{Full machine-readable cross-anchor grid.}
A complete observable-by-anchor cross-tabulation covering all
\crossanchorNObservables{} predicted observables (every entry in the
framework's structural catalogue --- 42 particle-mass entries plus
electroweak couplings, mixing angles, CP phases, gauge bosons,
lifetimes, decay widths, Compton wavelengths, atomic constants,
Yukawa couplings, dark-matter candidate spectrum, and gravitational
sector) is provided as
\href{https://github.com/gable-digital/papers/blob/75e7d76/substrate_falsifiable_predictions/cross_anchor_full.csv}{\texttt{cross\_anchor\_full.csv}}
(permanent reference, pinned to commit
\texttt{75e7d76} of \texttt{gable-digital/papers}, 2026-05-05).
The CSV exposes the seven anchors listed above (including the
parameter-free $v_C^{\rm struct}$ column not shown in
Table~\ref{tab:cross-anchor}), and for each observable carries a
\texttt{max\_pairwise\_anchor\_spread\_ppm} column so the rigidity
property is visible directly in the data without recomputation. All
values are computed at $\mathrm{mp.prec}=500$. The latest (untagged)
revision tracks the
\href{https://github.com/gable-digital/papers/tree/main/substrate_falsifiable_predictions}{\texttt{main}}
branch.

\section{SM mass spectrum, electroweak sector, mixing}

\begin{widetext}
% Auto-generated by substrate_mass_spectrum.py (rich longtable)
% Triple-anchor format with auto-flag colours.
{\scriptsize
\setlength{\LTcapwidth}{6.5in}
\begin{longtable}{lcccc}
\caption{Substrate-framework predictions for all derivable Standard Model and gravitational observables, computed at $\mathrm{mp.prec}=500$ from three independent anchor scales: CODATA electron mass, CODATA muon mass, and the directly-computed structural Higgs vev $v_C = (\dim - 2)/(1 - \delta_W)\cdot(1 + 5\alpha_\tau^3/8)$. Each cell carries the predicted value (15-digit precision) with the deviation against the observed value beneath, in the same cell. Auto-flag legend: \textcolor{OliveGreen}{$\checkmark$ green} ($|\mathrm{dev}| \leq 1\sigma_{\mathrm{obs}}$), \textcolor{orange}{$\bullet$ orange} ($1\sigma < |\mathrm{dev}| \leq 3\sigma$), \textcolor{red}{$\times$ red} ($|\mathrm{dev}| > 3\sigma$), \textcolor{blue}{$?$ blue} (estimate-pending; the prediction depends on upstream prediction work that is not yet closed, so the deviation is not yet falsifiable), \textcolor{violet}{$?$ violet} (mass predicted, but the framework cannot yet decide whether this candidate slot is physically occupied; used for half-integer-slot dark-mode candidates). Observed column carries the published experimental uncertainty ($\pm$ ppm/ppb/ppt) below each value where available. Anchor variance (max pairwise $\sim 158$~ppt between anchors across the full 42-particle catalogue, carried by the e-anchor-vs-vev-anchor pair; e-anchor-vs-$\mu$-anchor agrees at $\sim 4$~ppt) is exposed rather than hidden.}
\label{tab:complete-table} \\
\toprule
Observable & Observed & electron-anchor & muon-anchor & vev-anchor \\
\midrule
\endfirsthead
\multicolumn{5}{l}{\emph{Table~\ref{tab:complete-table} continued}} \\
\toprule
Observable & Observed & electron-anchor & muon-anchor & vev-anchor \\
\midrule
\endhead
\midrule
\multicolumn{5}{r}{\emph{(continued on next page)}} \\
\endfoot
\bottomrule
\endlastfoot
\multicolumn{5}{l}{\textit{CHARGED LEPTON MASSES}} \\
\parbox[t]{1.7in}{\raggedright electron (MeV)} & \shortstack[c]{0.5109989500000 \\ \scriptsize$\pm 300.000~\mathrm{ppt}$} & \shortstack[c]{0.5109989500000 \\ \scriptsize\textcolor{OliveGreen}{$\checkmark\,\sim 0\;(\mathrm{FP\ floor})$}} & \shortstack[c]{0.5109989500022 \\ \scriptsize\textcolor{OliveGreen}{$\checkmark\,+4.255~\mathrm{ppt}$}} & \shortstack[c]{0.5109989500806 \\ \scriptsize\textcolor{OliveGreen}{$\checkmark\,+157.642~\mathrm{ppt}$}} \\
\parbox[t]{1.7in}{\raggedright muon (MeV)} & \shortstack[c]{105.6583755000 \\ \scriptsize$\pm 22.000~\mathrm{ppb}$} & \shortstack[c]{105.6583754996 \\ \scriptsize\textcolor{OliveGreen}{$\checkmark\,-4.255~\mathrm{ppt}$}} & \shortstack[c]{105.6583755000 \\ \scriptsize\textcolor{OliveGreen}{$\checkmark\,\sim 0\;(\mathrm{FP\ floor})$}} & \shortstack[c]{105.6583755162 \\ \scriptsize\textcolor{OliveGreen}{$\checkmark\,+153.386~\mathrm{ppt}$}} \\
\parbox[t]{1.7in}{\raggedright tau (MeV)} & \shortstack[c]{1776.860000000 \\ \scriptsize$\pm 68.000~\mathrm{ppm}$} & \shortstack[c]{1776.859738841 \\ \scriptsize\textcolor{OliveGreen}{$\checkmark\,-146.978~\mathrm{ppb}$}} & \shortstack[c]{1776.859738849 \\ \scriptsize\textcolor{OliveGreen}{$\checkmark\,-146.973~\mathrm{ppb}$}} & \shortstack[c]{1776.859739121 \\ \scriptsize\textcolor{OliveGreen}{$\checkmark\,-146.820~\mathrm{ppb}$}} \\
\multicolumn{5}{l}{\textit{QUARK MASSES}} \\
\parbox[t]{1.7in}{\raggedright up (MeV)} & \shortstack[c]{2.160000000000 \\ \scriptsize$\pm 170000.000~\mathrm{ppm}$} & \shortstack[c]{2.172135018407 \\ \scriptsize\textcolor{OliveGreen}{$\checkmark\,+5618.064~\mathrm{ppm}$}} & \shortstack[c]{2.172135018416 \\ \scriptsize\textcolor{OliveGreen}{$\checkmark\,+5618.064~\mathrm{ppm}$}} & \shortstack[c]{2.172135018749 \\ \scriptsize\textcolor{OliveGreen}{$\checkmark\,+5618.064~\mathrm{ppm}$}} \\
\parbox[t]{1.7in}{\raggedright down (MeV)} & \shortstack[c]{4.670000000000 \\ \scriptsize$\pm 70000.000~\mathrm{ppm}$} & \shortstack[c]{4.666472667699 \\ \scriptsize\textcolor{OliveGreen}{$\checkmark\,-755.317~\mathrm{ppm}$}} & \shortstack[c]{4.666472667719 \\ \scriptsize\textcolor{OliveGreen}{$\checkmark\,-755.317~\mathrm{ppm}$}} & \shortstack[c]{4.666472668435 \\ \scriptsize\textcolor{OliveGreen}{$\checkmark\,-755.317~\mathrm{ppm}$}} \\
\parbox[t]{1.7in}{\raggedright strange (MeV)} & \shortstack[c]{93.40000000000 \\ \scriptsize$\pm 92000.000~\mathrm{ppm}$} & \shortstack[c]{93.25272875250 \\ \scriptsize\textcolor{OliveGreen}{$\checkmark\,-1576.780~\mathrm{ppm}$}} & \shortstack[c]{93.25272875289 \\ \scriptsize\textcolor{OliveGreen}{$\checkmark\,-1576.780~\mathrm{ppm}$}} & \shortstack[c]{93.25272876720 \\ \scriptsize\textcolor{OliveGreen}{$\checkmark\,-1576.780~\mathrm{ppm}$}} \\
\parbox[t]{1.7in}{\raggedright charm (MeV)} & \shortstack[c]{1273.000000000 \\ \scriptsize$\pm 3600.000~\mathrm{ppm}$} & \shortstack[c]{1274.297660596 \\ \scriptsize\textcolor{OliveGreen}{$\checkmark\,+1019.372~\mathrm{ppm}$}} & \shortstack[c]{1274.297660601 \\ \scriptsize\textcolor{OliveGreen}{$\checkmark\,+1019.372~\mathrm{ppm}$}} & \shortstack[c]{1274.297660796 \\ \scriptsize\textcolor{OliveGreen}{$\checkmark\,+1019.372~\mathrm{ppm}$}} \\
\parbox[t]{1.7in}{\raggedright bottom (MeV)} & \shortstack[c]{4183.000000000 \\ \scriptsize$\pm 1700.000~\mathrm{ppm}$} & \shortstack[c]{4178.924486322 \\ \scriptsize\textcolor{OliveGreen}{$\checkmark\,-974.304~\mathrm{ppm}$}} & \shortstack[c]{4178.924486340 \\ \scriptsize\textcolor{OliveGreen}{$\checkmark\,-974.304~\mathrm{ppm}$}} & \shortstack[c]{4178.924486981 \\ \scriptsize\textcolor{OliveGreen}{$\checkmark\,-974.304~\mathrm{ppm}$}} \\
\parbox[t]{1.7in}{\raggedright top (GeV)} & \shortstack[c]{172.5700000000 \\ \scriptsize$\pm 1700.000~\mathrm{ppm}$} & \shortstack[c]{172.5698832743 \\ \scriptsize\textcolor{OliveGreen}{$\checkmark\,-676.396~\mathrm{ppb}$}} & \shortstack[c]{172.5698832751 \\ \scriptsize\textcolor{OliveGreen}{$\checkmark\,-676.392~\mathrm{ppb}$}} & \shortstack[c]{172.5698833015 \\ \scriptsize\textcolor{OliveGreen}{$\checkmark\,-676.238~\mathrm{ppb}$}} \\
\multicolumn{5}{l}{\textit{NEUTRINO MASSES}} \\
\parbox[t]{1.7in}{\raggedright neutrino 1 (meV)} & \shortstack[c]{4.510000000000e-10 \\ \scriptsize$\pm 50000.000~\mathrm{ppm}$} & \shortstack[c]{4.516090890957e-10 \\ \scriptsize\textcolor{OliveGreen}{$\checkmark\,+1350.530~\mathrm{ppm}$}} & \shortstack[c]{4.516090890976e-10 \\ \scriptsize\textcolor{OliveGreen}{$\checkmark\,+1350.530~\mathrm{ppm}$}} & \shortstack[c]{4.516090891669e-10 \\ \scriptsize\textcolor{OliveGreen}{$\checkmark\,+1350.530~\mathrm{ppm}$}} \\
\parbox[t]{1.7in}{\raggedright neutrino 2 (meV)} & \shortstack[c]{8.613900000000e-9 \\ \scriptsize$\pm 50000.000~\mathrm{ppm}$} & \shortstack[c]{8.625379528706e-9 \\ \scriptsize\textcolor{OliveGreen}{$\checkmark\,+1332.675~\mathrm{ppm}$}} & \shortstack[c]{8.625379528743e-9 \\ \scriptsize\textcolor{OliveGreen}{$\checkmark\,+1332.675~\mathrm{ppm}$}} & \shortstack[c]{8.625379530066e-9 \\ \scriptsize\textcolor{OliveGreen}{$\checkmark\,+1332.675~\mathrm{ppm}$}} \\
\parbox[t]{1.7in}{\raggedright neutrino 3 (meV)} & \shortstack[c]{5.014980000000e-8 \\ \scriptsize$\pm 50000.000~\mathrm{ppm}$} & \shortstack[c]{5.021285707439e-8 \\ \scriptsize\textcolor{OliveGreen}{$\checkmark\,+1257.374~\mathrm{ppm}$}} & \shortstack[c]{5.021285707461e-8 \\ \scriptsize\textcolor{OliveGreen}{$\checkmark\,+1257.374~\mathrm{ppm}$}} & \shortstack[c]{5.021285708231e-8 \\ \scriptsize\textcolor{OliveGreen}{$\checkmark\,+1257.375~\mathrm{ppm}$}} \\
\multicolumn{5}{l}{\textit{ELECTROWEAK SECTOR}} \\
\parbox[t]{1.7in}{\raggedright Higgs vev (GeV)} & \shortstack[c]{246.2196500000 \\ \scriptsize$\pm 24.000~\mathrm{ppb}$} & \shortstack[c]{246.2196499809 \\ \scriptsize\textcolor{OliveGreen}{$\checkmark\,-77.617~\mathrm{ppt}$}} & \shortstack[c]{246.2196499819 \\ \scriptsize\textcolor{OliveGreen}{$\checkmark\,-73.361~\mathrm{ppt}$}} & \shortstack[c]{246.2196500197 \\ \scriptsize\textcolor{OliveGreen}{$\checkmark\,+80.025~\mathrm{ppt}$}} \\
\parbox[t]{1.7in}{\raggedright Higgs boson (GeV)} & \shortstack[c]{125.2000000000 \\ \scriptsize$\pm 880.000~\mathrm{ppm}$} & \shortstack[c]{125.1997813344 \\ \scriptsize\textcolor{OliveGreen}{$\checkmark\,-1.747~\mathrm{ppm}$}} & \shortstack[c]{125.1997813350 \\ \scriptsize\textcolor{OliveGreen}{$\checkmark\,-1.747~\mathrm{ppm}$}} & \shortstack[c]{125.1997813542 \\ \scriptsize\textcolor{OliveGreen}{$\checkmark\,-1.746~\mathrm{ppm}$}} \\
\parbox[t]{1.7in}{\raggedright W boson (GeV)} & \shortstack[c]{80.36920000000 \\ \scriptsize$\pm 160.000~\mathrm{ppm}$} & \shortstack[c]{80.36920037081 \\ \scriptsize\textcolor{OliveGreen}{$\checkmark\,+4.614~\mathrm{ppb}$}} & \shortstack[c]{80.36920037115 \\ \scriptsize\textcolor{OliveGreen}{$\checkmark\,+4.618~\mathrm{ppb}$}} & \shortstack[c]{80.36920038347 \\ \scriptsize\textcolor{OliveGreen}{$\checkmark\,+4.771~\mathrm{ppb}$}} \\
\parbox[t]{1.7in}{\raggedright Z boson (GeV)} & \shortstack[c]{91.18760000000 \\ \scriptsize$\pm 23.000~\mathrm{ppm}$} & \shortstack[c]{91.19003480815 \\ \scriptsize\textcolor{blue}{$?\,+26.701~\mathrm{ppm}$}} & \shortstack[c]{91.19003480853 \\ \scriptsize\textcolor{blue}{$?\,+26.701~\mathrm{ppm}$}} & \shortstack[c]{91.19003482252 \\ \scriptsize\textcolor{blue}{$?\,+26.701~\mathrm{ppm}$}} \\
\parbox[t]{1.7in}{\raggedright Weinberg mixing fraction} & \shortstack[c]{0.2232000000000 \\ \scriptsize$\pm 140.000~\mathrm{ppm}$} & \shortstack[c]{0.2232441517887 \\ \scriptsize\textcolor{blue}{$?\,+197.813~\mathrm{ppm}$}} & \shortstack[c]{0.2232441517887 \\ \scriptsize\textcolor{blue}{$?\,+197.813~\mathrm{ppm}$}} & \shortstack[c]{0.2232441517887 \\ \scriptsize\textcolor{blue}{$?\,+197.813~\mathrm{ppm}$}} \\
\parbox[t]{1.7in}{\raggedright SU(2) coupling} & \shortstack[c]{0.6533000000000 \\ \scriptsize$\pm 1500.000~\mathrm{ppm}$} & \shortstack[c]{0.6528252345176 \\ \scriptsize\textcolor{blue}{$?\,-726.719~\mathrm{ppm}$}} & \shortstack[c]{0.6528252345176 \\ \scriptsize\textcolor{blue}{$?\,-726.719~\mathrm{ppm}$}} & \shortstack[c]{0.6528252345176 \\ \scriptsize\textcolor{blue}{$?\,-726.719~\mathrm{ppm}$}} \\
\parbox[t]{1.7in}{\raggedright $U(1)_Y$ coupling (EW norm)} & \shortstack[c]{0.3500000000000 \\ \scriptsize$\pm 1000.000~\mathrm{ppm}$} & \shortstack[c]{0.3499811975896 \\ \scriptsize\textcolor{blue}{$?\,-53.721~\mathrm{ppm}$}} & \shortstack[c]{0.3499811975896 \\ \scriptsize\textcolor{blue}{$?\,-53.721~\mathrm{ppm}$}} & \shortstack[c]{0.3499811975896 \\ \scriptsize\textcolor{blue}{$?\,-53.721~\mathrm{ppm}$}} \\
\parbox[t]{1.7in}{\raggedright $U(1)_Y$ coupling (GUT norm)} & \shortstack[c]{0.4519000000000 \\ \scriptsize$\pm 1000.000~\mathrm{ppm}$} & \shortstack[c]{0.4518237832501 \\ \scriptsize\textcolor{blue}{$?\,-168.658~\mathrm{ppm}$}} & \shortstack[c]{0.4518237832501 \\ \scriptsize\textcolor{blue}{$?\,-168.658~\mathrm{ppm}$}} & \shortstack[c]{0.4518237832501 \\ \scriptsize\textcolor{blue}{$?\,-168.658~\mathrm{ppm}$}} \\
\parbox[t]{1.7in}{\raggedright SU(2) fine-structure} & \shortstack[c]{0.03394000000000 \\ \scriptsize$\pm 1500.000~\mathrm{ppm}$} & \shortstack[c]{0.03391438943683 \\ \scriptsize\textcolor{blue}{$?\,-754.583~\mathrm{ppm}$}} & \shortstack[c]{0.03391438943683 \\ \scriptsize\textcolor{blue}{$?\,-754.583~\mathrm{ppm}$}} & \shortstack[c]{0.03391438943683 \\ \scriptsize\textcolor{blue}{$?\,-754.583~\mathrm{ppm}$}} \\
\parbox[t]{1.7in}{\raggedright $U(1)_Y$ fine-structure (EW)} & \shortstack[c]{0.009745000000000 \\ \scriptsize$\pm 2000.000~\mathrm{ppm}$} & \shortstack[c]{0.009747192918717 \\ \scriptsize\textcolor{blue}{$?\,+225.030~\mathrm{ppm}$}} & \shortstack[c]{0.009747192918717 \\ \scriptsize\textcolor{blue}{$?\,+225.030~\mathrm{ppm}$}} & \shortstack[c]{0.009747192918717 \\ \scriptsize\textcolor{blue}{$?\,+225.030~\mathrm{ppm}$}} \\
\parbox[t]{1.7in}{\raggedright $U(1)_Y$ fine-structure (GUT)} & \shortstack[c]{0.01624100000000 \\ \scriptsize$\pm 2000.000~\mathrm{ppm}$} & \shortstack[c]{0.01624532153119 \\ \scriptsize\textcolor{blue}{$?\,+266.088~\mathrm{ppm}$}} & \shortstack[c]{0.01624532153119 \\ \scriptsize\textcolor{blue}{$?\,+266.088~\mathrm{ppm}$}} & \shortstack[c]{0.01624532153119 \\ \scriptsize\textcolor{blue}{$?\,+266.088~\mathrm{ppm}$}} \\
\parbox[t]{1.7in}{\raggedright EM fine-structure (tree)} & \shortstack[c]{0.007575408000000 \\ \scriptsize$\pm 1.000~\mathrm{ppb}$} & \shortstack[c]{0.007571189103257 \\ \scriptsize\textcolor{blue}{$?\,-556.920~\mathrm{ppm}$}} & \shortstack[c]{0.007571189103257 \\ \scriptsize\textcolor{blue}{$?\,-556.920~\mathrm{ppm}$}} & \shortstack[c]{0.007571189103257 \\ \scriptsize\textcolor{blue}{$?\,-556.920~\mathrm{ppm}$}} \\
\multicolumn{5}{l}{\textit{LEPTON YUKAWA COUPLINGS}} \\
\parbox[t]{1.7in}{\raggedright electron Yukawa} & \shortstack[c]{2.935028319017e-6 \\ \scriptsize$\pm 300.000~\mathrm{ppt}$} & \shortstack[c]{2.935028319245e-6 \\ \scriptsize\textcolor{OliveGreen}{$\checkmark\,+77.617~\mathrm{ppt}$}} & \shortstack[c]{2.935028319245e-6 \\ \scriptsize\textcolor{OliveGreen}{$\checkmark\,+77.617~\mathrm{ppt}$}} & \shortstack[c]{2.935028319245e-6 \\ \scriptsize\textcolor{OliveGreen}{$\checkmark\,+77.617~\mathrm{ppt}$}} \\
\parbox[t]{1.7in}{\raggedright muon Yukawa} & \shortstack[c]{0.0006068707660433 \\ \scriptsize$\pm 22.000~\mathrm{ppb}$} & \shortstack[c]{0.0006068707660878 \\ \scriptsize\textcolor{OliveGreen}{$\checkmark\,+73.361~\mathrm{ppt}$}} & \shortstack[c]{0.0006068707660878 \\ \scriptsize\textcolor{OliveGreen}{$\checkmark\,+73.361~\mathrm{ppt}$}} & \shortstack[c]{0.0006068707660878 \\ \scriptsize\textcolor{OliveGreen}{$\checkmark\,+73.361~\mathrm{ppt}$}} \\
\parbox[t]{1.7in}{\raggedright tau Yukawa} & \shortstack[c]{0.01020576347354 \\ \scriptsize$\pm 68.000~\mathrm{ppm}$} & \shortstack[c]{0.01020576197431 \\ \scriptsize\textcolor{OliveGreen}{$\checkmark\,-146.900~\mathrm{ppb}$}} & \shortstack[c]{0.01020576197431 \\ \scriptsize\textcolor{OliveGreen}{$\checkmark\,-146.900~\mathrm{ppb}$}} & \shortstack[c]{0.01020576197431 \\ \scriptsize\textcolor{OliveGreen}{$\checkmark\,-146.900~\mathrm{ppb}$}} \\
\multicolumn{5}{l}{\textit{QUARK YUKAWA COUPLINGS}} \\
\parbox[t]{1.7in}{\raggedright up Yukawa} & \shortstack[c]{1.240640742819e-5 \\ \scriptsize$\pm 170000.000~\mathrm{ppm}$} & \shortstack[c]{1.247610742106e-5 \\ \scriptsize\textcolor{OliveGreen}{$\checkmark\,+5618.064~\mathrm{ppm}$}} & \shortstack[c]{1.247610742106e-5 \\ \scriptsize\textcolor{OliveGreen}{$\checkmark\,+5618.064~\mathrm{ppm}$}} & \shortstack[c]{1.247610742106e-5 \\ \scriptsize\textcolor{OliveGreen}{$\checkmark\,+5618.064~\mathrm{ppm}$}} \\
\parbox[t]{1.7in}{\raggedright down Yukawa} & \shortstack[c]{2.682311235631e-5 \\ \scriptsize$\pm 70000.000~\mathrm{ppm}$} & \shortstack[c]{2.680285239466e-5 \\ \scriptsize\textcolor{OliveGreen}{$\checkmark\,-755.317~\mathrm{ppm}$}} & \shortstack[c]{2.680285239466e-5 \\ \scriptsize\textcolor{OliveGreen}{$\checkmark\,-755.317~\mathrm{ppm}$}} & \shortstack[c]{2.680285239466e-5 \\ \scriptsize\textcolor{OliveGreen}{$\checkmark\,-755.317~\mathrm{ppm}$}} \\
\parbox[t]{1.7in}{\raggedright strange Yukawa} & \shortstack[c]{0.0005364622471263 \\ \scriptsize$\pm 92000.000~\mathrm{ppm}$} & \shortstack[c]{0.0005356163642517 \\ \scriptsize\textcolor{OliveGreen}{$\checkmark\,-1576.780~\mathrm{ppm}$}} & \shortstack[c]{0.0005356163642517 \\ \scriptsize\textcolor{OliveGreen}{$\checkmark\,-1576.780~\mathrm{ppm}$}} & \shortstack[c]{0.0005356163642517 \\ \scriptsize\textcolor{OliveGreen}{$\checkmark\,-1576.780~\mathrm{ppm}$}} \\
\parbox[t]{1.7in}{\raggedright charm Yukawa} & \shortstack[c]{0.007311739192631 \\ \scriptsize$\pm 3600.000~\mathrm{ppm}$} & \shortstack[c]{0.007319192575631 \\ \scriptsize\textcolor{OliveGreen}{$\checkmark\,+1019.372~\mathrm{ppm}$}} & \shortstack[c]{0.007319192575631 \\ \scriptsize\textcolor{OliveGreen}{$\checkmark\,+1019.372~\mathrm{ppm}$}} & \shortstack[c]{0.007319192575631 \\ \scriptsize\textcolor{OliveGreen}{$\checkmark\,+1019.372~\mathrm{ppm}$}} \\
\parbox[t]{1.7in}{\raggedright bottom Yukawa} & \shortstack[c]{0.02402592697783 \\ \scriptsize$\pm 1700.000~\mathrm{ppm}$} & \shortstack[c]{0.02400251842267 \\ \scriptsize\textcolor{OliveGreen}{$\checkmark\,-974.304~\mathrm{ppm}$}} & \shortstack[c]{0.02400251842267 \\ \scriptsize\textcolor{OliveGreen}{$\checkmark\,-974.304~\mathrm{ppm}$}} & \shortstack[c]{0.02400251842267 \\ \scriptsize\textcolor{OliveGreen}{$\checkmark\,-974.304~\mathrm{ppm}$}} \\
\parbox[t]{1.7in}{\raggedright top Yukawa} & \shortstack[c]{0.9911915416122 \\ \scriptsize$\pm 1700.000~\mathrm{ppm}$} & \shortstack[c]{0.9911908712511 \\ \scriptsize\textcolor{OliveGreen}{$\checkmark\,-676.318~\mathrm{ppb}$}} & \shortstack[c]{0.9911908712511 \\ \scriptsize\textcolor{OliveGreen}{$\checkmark\,-676.318~\mathrm{ppb}$}} & \shortstack[c]{0.9911908712511 \\ \scriptsize\textcolor{OliveGreen}{$\checkmark\,-676.318~\mathrm{ppb}$}} \\
\multicolumn{5}{l}{\textit{HIGGS SELF-COUPLING}} \\
\parbox[t]{1.7in}{\raggedright Higgs self-coupling} & \shortstack[c]{0.1292805654113 \\ \scriptsize$\pm 1800.000~\mathrm{ppm}$} & \shortstack[c]{0.1292801138469 \\ \scriptsize\textcolor{OliveGreen}{$\checkmark\,-3.493~\mathrm{ppm}$}} & \shortstack[c]{0.1292801138469 \\ \scriptsize\textcolor{OliveGreen}{$\checkmark\,-3.493~\mathrm{ppm}$}} & \shortstack[c]{0.1292801138469 \\ \scriptsize\textcolor{OliveGreen}{$\checkmark\,-3.493~\mathrm{ppm}$}} \\
\multicolumn{5}{l}{\textit{STRONG COUPLING}} \\
\parbox[t]{1.7in}{\raggedright strong coupling} & \shortstack[c]{0.1180000000000 \\ \scriptsize$\pm 7600.000~\mathrm{ppm}$} & \shortstack[c]{0.1180478309232 \\ \scriptsize\textcolor{blue}{$?\,+405.347~\mathrm{ppm}$}} & \shortstack[c]{0.1180478309232 \\ \scriptsize\textcolor{blue}{$?\,+405.347~\mathrm{ppm}$}} & \shortstack[c]{0.1180478309232 \\ \scriptsize\textcolor{blue}{$?\,+405.347~\mathrm{ppm}$}} \\
\multicolumn{5}{l}{\textit{PMNS NEUTRINO MIXING}} \\
\parbox[t]{1.7in}{\raggedright PMNS solar mixing fraction} & \shortstack[c]{0.3030000000000 \\ \scriptsize$\pm 50000.000~\mathrm{ppm}$} & \shortstack[c]{0.3065788699420 \\ \scriptsize\textcolor{blue}{$?\,+11811.452~\mathrm{ppm}$}} & \shortstack[c]{0.3065788699420 \\ \scriptsize\textcolor{blue}{$?\,+11811.452~\mathrm{ppm}$}} & \shortstack[c]{0.3065788699420 \\ \scriptsize\textcolor{blue}{$?\,+11811.452~\mathrm{ppm}$}} \\
\parbox[t]{1.7in}{\raggedright PMNS atmospheric mixing fraction} & \shortstack[c]{0.5720000000000 \\ \scriptsize$\pm 30000.000~\mathrm{ppm}$} & \shortstack[c]{0.5714285714286 \\ \scriptsize\textcolor{blue}{$?\,-999.001~\mathrm{ppm}$}} & \shortstack[c]{0.5714285714286 \\ \scriptsize\textcolor{blue}{$?\,-999.001~\mathrm{ppm}$}} & \shortstack[c]{0.5714285714286 \\ \scriptsize\textcolor{blue}{$?\,-999.001~\mathrm{ppm}$}} \\
\parbox[t]{1.7in}{\raggedright PMNS reactor mixing fraction} & \shortstack[c]{0.02203000000000 \\ \scriptsize$\pm 30000.000~\mathrm{ppm}$} & \shortstack[c]{0.02197802197802 \\ \scriptsize\textcolor{blue}{$?\,-2359.420~\mathrm{ppm}$}} & \shortstack[c]{0.02197802197802 \\ \scriptsize\textcolor{blue}{$?\,-2359.420~\mathrm{ppm}$}} & \shortstack[c]{0.02197802197802 \\ \scriptsize\textcolor{blue}{$?\,-2359.420~\mathrm{ppm}$}} \\
\parbox[t]{1.7in}{\raggedright PMNS CP phase (rad)} & \shortstack[c]{3.438298626429 \\ \scriptsize$\pm 50000.000~\mathrm{ppm}$} & \shortstack[c]{3.455751918949 \\ \scriptsize\textcolor{blue}{$?\,+5076.142~\mathrm{ppm}$}} & \shortstack[c]{3.455751918949 \\ \scriptsize\textcolor{blue}{$?\,+5076.142~\mathrm{ppm}$}} & \shortstack[c]{3.455751918949 \\ \scriptsize\textcolor{blue}{$?\,+5076.142~\mathrm{ppm}$}} \\
\multicolumn{5}{l}{\textit{CKM QUARK MIXING}} \\
\parbox[t]{1.7in}{\raggedright Cabibbo mixing fraction} & \shortstack[c]{0.05064000000000 \\ \scriptsize$\pm 4000.000~\mathrm{ppm}$} & \shortstack[c]{0.05035971223022 \\ \scriptsize\textcolor{blue}{$?\,-5534.909~\mathrm{ppm}$}} & \shortstack[c]{0.05035971223022 \\ \scriptsize\textcolor{blue}{$?\,-5534.909~\mathrm{ppm}$}} & \shortstack[c]{0.05035971223022 \\ \scriptsize\textcolor{blue}{$?\,-5534.909~\mathrm{ppm}$}} \\
\parbox[t]{1.7in}{\raggedright CKM atmospheric mixing fraction} & \shortstack[c]{0.001682000000000 \\ \scriptsize$\pm 15000.000~\mathrm{ppm}$} & \shortstack[c]{0.001681796794019 \\ \scriptsize\textcolor{blue}{$?\,-120.812~\mathrm{ppm}$}} & \shortstack[c]{0.001681796794019 \\ \scriptsize\textcolor{blue}{$?\,-120.812~\mathrm{ppm}$}} & \shortstack[c]{0.001681796794019 \\ \scriptsize\textcolor{blue}{$?\,-120.812~\mathrm{ppm}$}} \\
\parbox[t]{1.7in}{\raggedright CKM reactor mixing fraction} & \shortstack[c]{1.460000000000e-5 \\ \scriptsize$\pm 30000.000~\mathrm{ppm}$} & \shortstack[c]{1.452605387890e-5 \\ \scriptsize\textcolor{blue}{$?\,-5064.803~\mathrm{ppm}$}} & \shortstack[c]{1.452605387890e-5 \\ \scriptsize\textcolor{blue}{$?\,-5064.803~\mathrm{ppm}$}} & \shortstack[c]{1.452605387890e-5 \\ \scriptsize\textcolor{blue}{$?\,-5064.803~\mathrm{ppm}$}} \\
\parbox[t]{1.7in}{\raggedright CKM CP phase (rad)} & \shortstack[c]{1.143190660056 \\ \scriptsize$\pm 50000.000~\mathrm{ppm}$} & \shortstack[c]{1.151917306316 \\ \scriptsize\textcolor{blue}{$?\,+7633.588~\mathrm{ppm}$}} & \shortstack[c]{1.151917306316 \\ \scriptsize\textcolor{blue}{$?\,+7633.588~\mathrm{ppm}$}} & \shortstack[c]{1.151917306316 \\ \scriptsize\textcolor{blue}{$?\,+7633.588~\mathrm{ppm}$}} \\
\multicolumn{5}{l}{\textit{QCD THETA ANGLE}} \\
\parbox[t]{1.7in}{\raggedright QCD theta angle (rad)} & 0.0 & 0.0 & 0.0 & 0.0 \\
\multicolumn{5}{l}{\textit{MASSLESS GAUGE BOSONS}} \\
\parbox[t]{1.7in}{\raggedright photon (MeV)} & 0.0 & 0.0 & 0.0 & 0.0 \\
\parbox[t]{1.7in}{\raggedright gluon 1 (MeV)} & 0.0 & 0.0 & 0.0 & 0.0 \\
\parbox[t]{1.7in}{\raggedright gluon 2 (MeV)} & 0.0 & 0.0 & 0.0 & 0.0 \\
\parbox[t]{1.7in}{\raggedright gluon 3 (MeV)} & 0.0 & 0.0 & 0.0 & 0.0 \\
\parbox[t]{1.7in}{\raggedright gluon 4 (MeV)} & 0.0 & 0.0 & 0.0 & 0.0 \\
\parbox[t]{1.7in}{\raggedright gluon 5 (MeV)} & 0.0 & 0.0 & 0.0 & 0.0 \\
\parbox[t]{1.7in}{\raggedright gluon 6 (MeV)} & 0.0 & 0.0 & 0.0 & 0.0 \\
\parbox[t]{1.7in}{\raggedright gluon 7 (MeV)} & 0.0 & 0.0 & 0.0 & 0.0 \\
\parbox[t]{1.7in}{\raggedright gluon 8 (MeV)} & 0.0 & 0.0 & 0.0 & 0.0 \\
\multicolumn{5}{l}{\textit{GRAVITATIONAL SECTOR}} \\
\parbox[t]{1.7in}{\raggedright Newton's G (m$^3$kg$^{-1}$s$^{-2}$)} & \shortstack[c]{6.674300000000e-11 \\ \scriptsize$\pm 22.000~\mathrm{ppm}$} & \shortstack[c]{6.674316910358e-11 \\ \scriptsize\textcolor{OliveGreen}{$\checkmark\,+2.534~\mathrm{ppm}$}} & \shortstack[c]{6.674316910301e-11 \\ \scriptsize\textcolor{OliveGreen}{$\checkmark\,+2.534~\mathrm{ppm}$}} & \shortstack[c]{6.674316908253e-11 \\ \scriptsize\textcolor{OliveGreen}{$\checkmark\,+2.533~\mathrm{ppm}$}} \\
\parbox[t]{1.7in}{\raggedright Planck mass (GeV)} & \shortstack[c]{1.220890000000e+22 \\ \scriptsize$\pm 11.000~\mathrm{ppm}$} & \shortstack[c]{1.220888581931e+22 \\ \scriptsize\textcolor{OliveGreen}{$\checkmark\,-1.162~\mathrm{ppm}$}} & \shortstack[c]{1.220888581936e+22 \\ \scriptsize\textcolor{OliveGreen}{$\checkmark\,-1.161~\mathrm{ppm}$}} & \shortstack[c]{1.220888582124e+22 \\ \scriptsize\textcolor{OliveGreen}{$\checkmark\,-1.161~\mathrm{ppm}$}} \\
\multicolumn{5}{l}{\textit{ANOMALOUS MAGNETIC MOMENTS}} \\
\parbox[t]{1.7in}{\raggedright electron anomalous moment} & \shortstack[c]{0.001159652181000 \\ \scriptsize$\pm 1.700~\mathrm{ppt}$} & \shortstack[c]{0.001160667181068 \\ \scriptsize\textcolor{blue}{$?\,+875.263~\mathrm{ppm}$}} & \shortstack[c]{0.001160667181068 \\ \scriptsize\textcolor{blue}{$?\,+875.263~\mathrm{ppm}$}} & \shortstack[c]{0.001160667181068 \\ \scriptsize\textcolor{blue}{$?\,+875.263~\mathrm{ppm}$}} \\
\parbox[t]{1.7in}{\raggedright muon anomalous moment} & \shortstack[c]{0.001165920620000 \\ \scriptsize$\pm 460.000~\mathrm{ppb}$} & \shortstack[c]{0.001166944225438 \\ \scriptsize\textcolor{blue}{$?\,+877.938~\mathrm{ppm}$}} & \shortstack[c]{0.001166944225438 \\ \scriptsize\textcolor{blue}{$?\,+877.938~\mathrm{ppm}$}} & \shortstack[c]{0.001166944225438 \\ \scriptsize\textcolor{blue}{$?\,+877.938~\mathrm{ppm}$}} \\
\parbox[t]{1.7in}{\raggedright tau anomalous moment} & \shortstack[c]{0.001177210000000 \\ \scriptsize$\pm 2100.000~\mathrm{ppm}$} & \shortstack[c]{0.001176416612358 \\ \scriptsize\textcolor{blue}{$?\,-673.956~\mathrm{ppm}$}} & \shortstack[c]{0.001176416612358 \\ \scriptsize\textcolor{blue}{$?\,-673.956~\mathrm{ppm}$}} & \shortstack[c]{0.001176416612358 \\ \scriptsize\textcolor{blue}{$?\,-673.956~\mathrm{ppm}$}} \\
\multicolumn{5}{l}{\textit{FERMI CONSTANT}} \\
\parbox[t]{1.7in}{\raggedright Fermi constant (GeV$^{-2}$)} & \shortstack[c]{1.166378700000e-5 \\ \scriptsize$\pm 500.000~\mathrm{ppb}$} & \shortstack[c]{1.166378707705e-5 \\ \scriptsize\textcolor{OliveGreen}{$\checkmark\,+6.606~\mathrm{ppb}$}} & \shortstack[c]{1.166378707695e-5 \\ \scriptsize\textcolor{OliveGreen}{$\checkmark\,+6.597~\mathrm{ppb}$}} & \shortstack[c]{1.166378707337e-5 \\ \scriptsize\textcolor{OliveGreen}{$\checkmark\,+6.291~\mathrm{ppb}$}} \\
\multicolumn{5}{l}{\textit{LEPTON LIFETIMES}} \\
\parbox[t]{1.7in}{\raggedright muon lifetime (s)} & \shortstack[c]{2.196981100000e-6 \\ \scriptsize$\pm 1.000~\mathrm{ppm}$} & \shortstack[c]{2.197000758647e-6 \\ \scriptsize\textcolor{blue}{$?\,+8.948~\mathrm{ppm}$}} & \shortstack[c]{2.197000758638e-6 \\ \scriptsize\textcolor{blue}{$?\,+8.948~\mathrm{ppm}$}} & \shortstack[c]{2.197000758301e-6 \\ \scriptsize\textcolor{blue}{$?\,+8.948~\mathrm{ppm}$}} \\
\parbox[t]{1.7in}{\raggedright tau lifetime (s)} & \shortstack[c]{2.903000000000e-13 \\ \scriptsize$\pm 400.000~\mathrm{ppm}$} & \shortstack[c]{2.975066653724e-13 \\ \scriptsize\textcolor{blue}{$?\,+24824.889~\mathrm{ppm}$}} & \shortstack[c]{2.975066653712e-13 \\ \scriptsize\textcolor{blue}{$?\,+24824.889~\mathrm{ppm}$}} & \shortstack[c]{2.975066653255e-13 \\ \scriptsize\textcolor{blue}{$?\,+24824.889~\mathrm{ppm}$}} \\
\multicolumn{5}{l}{\textit{W AND Z DECAY WIDTHS}} \\
\parbox[t]{1.7in}{\raggedright W total decay width (GeV)} & \shortstack[c]{2.085000000000 \\ \scriptsize$\pm 2000.000~\mathrm{ppm}$} & \shortstack[c]{2.097035187050 \\ \scriptsize\textcolor{blue}{$?\,+5772.272~\mathrm{ppm}$}} & \shortstack[c]{2.097035187059 \\ \scriptsize\textcolor{blue}{$?\,+5772.272~\mathrm{ppm}$}} & \shortstack[c]{2.097035187381 \\ \scriptsize\textcolor{blue}{$?\,+5772.272~\mathrm{ppm}$}} \\
\parbox[t]{1.7in}{\raggedright Z total decay width (GeV)} & \shortstack[c]{2.495500000000 \\ \scriptsize$\pm 90.000~\mathrm{ppm}$} & \shortstack[c]{2.508962908085 \\ \scriptsize\textcolor{blue}{$?\,+5394.874~\mathrm{ppm}$}} & \shortstack[c]{2.508962908096 \\ \scriptsize\textcolor{blue}{$?\,+5394.874~\mathrm{ppm}$}} & \shortstack[c]{2.508962908481 \\ \scriptsize\textcolor{blue}{$?\,+5394.874~\mathrm{ppm}$}} \\
\parbox[t]{1.7in}{\raggedright Z invisible width (3 neutrinos) (GeV)} & \shortstack[c]{0.4995000000000 \\ \scriptsize$\pm 3000.000~\mathrm{ppm}$} & \shortstack[c]{0.4976876518476 \\ \scriptsize\textcolor{blue}{$?\,-3628.325~\mathrm{ppm}$}} & \shortstack[c]{0.4976876518498 \\ \scriptsize\textcolor{blue}{$?\,-3628.325~\mathrm{ppm}$}} & \shortstack[c]{0.4976876519261 \\ \scriptsize\textcolor{blue}{$?\,-3628.324~\mathrm{ppm}$}} \\
\parbox[t]{1.7in}{\raggedright Z leptonic width (per lepton) (GeV)} & \shortstack[c]{0.08398000000000 \\ \scriptsize$\pm 500.000~\mathrm{ppm}$} & \shortstack[c]{0.08389802825078 \\ \scriptsize\textcolor{blue}{$?\,-976.087~\mathrm{ppm}$}} & \shortstack[c]{0.08389802825113 \\ \scriptsize\textcolor{blue}{$?\,-976.087~\mathrm{ppm}$}} & \shortstack[c]{0.08389802826400 \\ \scriptsize\textcolor{blue}{$?\,-976.086~\mathrm{ppm}$}} \\
\multicolumn{5}{l}{\textit{HIGGS DECAY WIDTHS}} \\
\parbox[t]{1.7in}{\raggedright Higgs width to bb (GeV)} & \shortstack[c]{0.002200000000000 \\ \scriptsize$\pm 140000.000~\mathrm{ppm}$} & \shortstack[c]{0.002303033296527 \\ \scriptsize\textcolor{blue}{$?\,+46833.317~\mathrm{ppm}$}} & \shortstack[c]{0.002303033296517 \\ \scriptsize\textcolor{blue}{$?\,+46833.317~\mathrm{ppm}$}} & \shortstack[c]{0.002303033296166 \\ \scriptsize\textcolor{blue}{$?\,+46833.316~\mathrm{ppm}$}} \\
\parbox[t]{1.7in}{\raggedright Higgs width to tau tau (GeV)} & \shortstack[c]{0.0002300000000000 \\ \scriptsize$\pm 100000.000~\mathrm{ppm}$} & \shortstack[c]{0.0002591191584134 \\ \scriptsize\textcolor{blue}{$?\,+126605.037~\mathrm{ppm}$}} & \shortstack[c]{0.0002591191584145 \\ \scriptsize\textcolor{blue}{$?\,+126605.037~\mathrm{ppm}$}} & \shortstack[c]{0.0002591191584543 \\ \scriptsize\textcolor{blue}{$?\,+126605.037~\mathrm{ppm}$}} \\
\parbox[t]{1.7in}{\raggedright Higgs width to cc (GeV)} & \shortstack[c]{0.0001200000000000 \\ \scriptsize$\pm 300000.000~\mathrm{ppm}$} & \shortstack[c]{0.0001140529647702 \\ \scriptsize\textcolor{blue}{$?\,-49558.627~\mathrm{ppm}$}} & \shortstack[c]{0.0001140529647697 \\ \scriptsize\textcolor{blue}{$?\,-49558.627~\mathrm{ppm}$}} & \shortstack[c]{0.0001140529647522 \\ \scriptsize\textcolor{blue}{$?\,-49558.627~\mathrm{ppm}$}} \\
\parbox[t]{1.7in}{\raggedright Higgs total fermionic width (GeV)} & \shortstack[c]{0.002510000000000 \\ \scriptsize$\pm 50000.000~\mathrm{ppm}$} & \shortstack[c]{0.002676205419710 \\ \scriptsize\textcolor{blue}{$?\,+66217.299~\mathrm{ppm}$}} & \shortstack[c]{0.002676205419701 \\ \scriptsize\textcolor{blue}{$?\,+66217.299~\mathrm{ppm}$}} & \shortstack[c]{0.002676205419372 \\ \scriptsize\textcolor{blue}{$?\,+66217.299~\mathrm{ppm}$}} \\
\multicolumn{5}{l}{\textit{COMPTON WAVELENGTHS AND ATOMIC CONSTANTS}} \\
\parbox[t]{1.7in}{\raggedright electron Compton wavelength (fm)} & \shortstack[c]{2426.310238670 \\ \scriptsize$\pm 1.500~\mathrm{ppb}$} & \shortstack[c]{2426.310237936 \\ \scriptsize\textcolor{OliveGreen}{$\checkmark\,-302.639~\mathrm{ppt}$}} & \shortstack[c]{2426.310237925 \\ \scriptsize\textcolor{OliveGreen}{$\checkmark\,-306.895~\mathrm{ppt}$}} & \shortstack[c]{2426.310237553 \\ \scriptsize\textcolor{OliveGreen}{$\checkmark\,-460.281~\mathrm{ppt}$}} \\
\parbox[t]{1.7in}{\raggedright muon Compton wavelength (fm)} & \shortstack[c]{11.73444110000 \\ \scriptsize$\pm 22.000~\mathrm{ppb}$} & \shortstack[c]{11.73444109942 \\ \scriptsize\textcolor{OliveGreen}{$\checkmark\,-49.815~\mathrm{ppt}$}} & \shortstack[c]{11.73444109937 \\ \scriptsize\textcolor{OliveGreen}{$\checkmark\,-54.070~\mathrm{ppt}$}} & \shortstack[c]{11.73444109757 \\ \scriptsize\textcolor{OliveGreen}{$\checkmark\,-207.457~\mathrm{ppt}$}} \\
\parbox[t]{1.7in}{\raggedright tau Compton wavelength (fm)} & \shortstack[c]{0.6977713400000 \\ \scriptsize$\pm 70.000~\mathrm{ppm}$} & \shortstack[c]{0.6977714429885 \\ \scriptsize\textcolor{OliveGreen}{$\checkmark\,+147.596~\mathrm{ppb}$}} & \shortstack[c]{0.6977714429855 \\ \scriptsize\textcolor{OliveGreen}{$\checkmark\,+147.592~\mathrm{ppb}$}} & \shortstack[c]{0.6977714428785 \\ \scriptsize\textcolor{OliveGreen}{$\checkmark\,+147.439~\mathrm{ppb}$}} \\
\parbox[t]{1.7in}{\raggedright Bohr radius (fm)} & \shortstack[c]{52917.72109000 \\ \scriptsize$\pm 1.500~\mathrm{ppb}$} & \shortstack[c]{52871.37492741 \\ \scriptsize\textcolor{blue}{$?\,-875.816~\mathrm{ppm}$}} & \shortstack[c]{52871.37492719 \\ \scriptsize\textcolor{blue}{$?\,-875.816~\mathrm{ppm}$}} & \shortstack[c]{52871.37491908 \\ \scriptsize\textcolor{blue}{$?\,-875.816~\mathrm{ppm}$}} \\
\parbox[t]{1.7in}{\raggedright Rydberg energy (eV)} & \shortstack[c]{13.60569312300 \\ \scriptsize$\pm 14.000~\mathrm{ppt}$} & \shortstack[c]{13.62955661536 \\ \scriptsize\textcolor{blue}{$?\,+1753.934~\mathrm{ppm}$}} & \shortstack[c]{13.62955661542 \\ \scriptsize\textcolor{blue}{$?\,+1753.934~\mathrm{ppm}$}} & \shortstack[c]{13.62955661751 \\ \scriptsize\textcolor{blue}{$?\,+1753.934~\mathrm{ppm}$}} \\
\multicolumn{5}{l}{\textit{EFFECTIVE NEUTRINO SPECIES}} \\
\parbox[t]{1.7in}{\raggedright effective neutrino species} & \shortstack[c]{3.046000000000 \\ \scriptsize$\pm 56000.000~\mathrm{ppm}$} & \shortstack[c]{3.000000000000 \\ \scriptsize\textcolor{blue}{$?\,-15101.773~\mathrm{ppm}$}} & \shortstack[c]{3.000000000000 \\ \scriptsize\textcolor{blue}{$?\,-15101.773~\mathrm{ppm}$}} & \shortstack[c]{3.000000000000 \\ \scriptsize\textcolor{blue}{$?\,-15101.773~\mathrm{ppm}$}} \\
\multicolumn{5}{l}{\textit{DARK MATTER ULTRALIGHT BAND (predicted-but-unobserved; falsification targets)}} \\
\parbox[t]{1.7in}{\raggedright DM ultralight 1 (heaviest) (eV)} & -- & \shortstack[c]{6.205011597344e-7 \\ \scriptsize\textcolor{blue}{$?\,\mathrm{predicted}$}} & \shortstack[c]{6.205011597370e-7 \\ \scriptsize\textcolor{blue}{$?\,\mathrm{predicted}$}} & \shortstack[c]{6.205011598322e-7 \\ \scriptsize\textcolor{blue}{$?\,\mathrm{predicted}$}} \\
\parbox[t]{1.7in}{\raggedright DM ultralight 2 (eV)} & -- & \shortstack[c]{2.232018560196e-9 \\ \scriptsize\textcolor{blue}{$?\,\mathrm{predicted}$}} & \shortstack[c]{2.232018560205e-9 \\ \scriptsize\textcolor{blue}{$?\,\mathrm{predicted}$}} & \shortstack[c]{2.232018560548e-9 \\ \scriptsize\textcolor{blue}{$?\,\mathrm{predicted}$}} \\
\parbox[t]{1.7in}{\raggedright DM ultralight 3 (eV)} & -- & \shortstack[c]{8.028843741711e-12 \\ \scriptsize\textcolor{blue}{$?\,\mathrm{predicted}$}} & \shortstack[c]{8.028843741745e-12 \\ \scriptsize\textcolor{blue}{$?\,\mathrm{predicted}$}} & \shortstack[c]{8.028843742977e-12 \\ \scriptsize\textcolor{blue}{$?\,\mathrm{predicted}$}} \\
\parbox[t]{1.7in}{\raggedright DM ultralight 4 (eV)} & -- & \shortstack[c]{2.888073288385e-14 \\ \scriptsize\textcolor{blue}{$?\,\mathrm{predicted}$}} & \shortstack[c]{2.888073288398e-14 \\ \scriptsize\textcolor{blue}{$?\,\mathrm{predicted}$}} & \shortstack[c]{2.888073288841e-14 \\ \scriptsize\textcolor{blue}{$?\,\mathrm{predicted}$}} \\
\parbox[t]{1.7in}{\raggedright DM ultralight 5 (eV)} & -- & \shortstack[c]{1.038875283592e-16 \\ \scriptsize\textcolor{blue}{$?\,\mathrm{predicted}$}} & \shortstack[c]{1.038875283596e-16 \\ \scriptsize\textcolor{blue}{$?\,\mathrm{predicted}$}} & \shortstack[c]{1.038875283756e-16 \\ \scriptsize\textcolor{blue}{$?\,\mathrm{predicted}$}} \\
\parbox[t]{1.7in}{\raggedright DM ultralight 6 (eV)} & -- & \shortstack[c]{3.736961451769e-19 \\ \scriptsize\textcolor{blue}{$?\,\mathrm{predicted}$}} & \shortstack[c]{3.736961451785e-19 \\ \scriptsize\textcolor{blue}{$?\,\mathrm{predicted}$}} & \shortstack[c]{3.736961452358e-19 \\ \scriptsize\textcolor{blue}{$?\,\mathrm{predicted}$}} \\
\parbox[t]{1.7in}{\raggedright DM ultralight 7 (lightest) (eV)} & -- & \shortstack[c]{1.344230738046e-21 \\ \scriptsize\textcolor{blue}{$?\,\mathrm{predicted}$}} & \shortstack[c]{1.344230738052e-21 \\ \scriptsize\textcolor{blue}{$?\,\mathrm{predicted}$}} & \shortstack[c]{1.344230738258e-21 \\ \scriptsize\textcolor{blue}{$?\,\mathrm{predicted}$}} \\
\multicolumn{5}{l}{\textit{DARK MATTER MeV--GeV BAND (predicted-but-unobserved; direct-detection targets)}} \\
\parbox[t]{1.7in}{\raggedright DM MeV--GeV mode A (MeV)} & -- & \shortstack[c]{0.8335357316677 \\ \scriptsize\textcolor{blue}{$?\,\mathrm{predicted}$}} & \shortstack[c]{0.8335357316712 \\ \scriptsize\textcolor{blue}{$?\,\mathrm{predicted}$}} & \shortstack[c]{0.8335357317991 \\ \scriptsize\textcolor{blue}{$?\,\mathrm{predicted}$}} \\
\parbox[t]{1.7in}{\raggedright DM MeV--GeV mode B (MeV)} & -- & \shortstack[c]{5.828295547031 \\ \scriptsize\textcolor{blue}{$?\,\mathrm{predicted}$}} & \shortstack[c]{5.828295547056 \\ \scriptsize\textcolor{blue}{$?\,\mathrm{predicted}$}} & \shortstack[c]{5.828295547950 \\ \scriptsize\textcolor{blue}{$?\,\mathrm{predicted}$}} \\
\parbox[t]{1.7in}{\raggedright DM MeV--GeV mode C (MeV)} & -- & \shortstack[c]{46.62636437625 \\ \scriptsize\textcolor{blue}{$?\,\mathrm{predicted}$}} & \shortstack[c]{46.62636437645 \\ \scriptsize\textcolor{blue}{$?\,\mathrm{predicted}$}} & \shortstack[c]{46.62636438360 \\ \scriptsize\textcolor{blue}{$?\,\mathrm{predicted}$}} \\
\parbox[t]{1.7in}{\raggedright DM MeV--GeV mode D (GeV)} & -- & \shortstack[c]{1.492043660040 \\ \scriptsize\textcolor{blue}{$?\,\mathrm{predicted}$}} & \shortstack[c]{1.492043660046 \\ \scriptsize\textcolor{blue}{$?\,\mathrm{predicted}$}} & \shortstack[c]{1.492043660275 \\ \scriptsize\textcolor{blue}{$?\,\mathrm{predicted}$}} \\
\parbox[t]{1.7in}{\raggedright DM MeV--GeV mode E (GeV)} & -- & \shortstack[c]{4.816169313316 \\ \scriptsize\textcolor{blue}{$?\,\mathrm{predicted}$}} & \shortstack[c]{4.816169313336 \\ \scriptsize\textcolor{blue}{$?\,\mathrm{predicted}$}} & \shortstack[c]{4.816169314075 \\ \scriptsize\textcolor{blue}{$?\,\mathrm{predicted}$}} \\
\parbox[t]{1.7in}{\raggedright DM MeV--GeV mode F (GeV)} & -- & \shortstack[c]{592.3496671013 \\ \scriptsize\textcolor{blue}{$?\,\mathrm{predicted}$}} & \shortstack[c]{592.3496671038 \\ \scriptsize\textcolor{blue}{$?\,\mathrm{predicted}$}} & \shortstack[c]{592.3496671946 \\ \scriptsize\textcolor{blue}{$?\,\mathrm{predicted}$}} \\
\multicolumn{5}{l}{\textit{HALF-SLOT CANDIDATES (spin-1/2 dark-mode candidates; slot occupancy open)}} \\
\parbox[t]{1.7in}{\raggedright DM half-slot candidate 1 (MeV)} & -- & \shortstack[c]{0.6552852638606 \\ \scriptsize\textcolor{violet}{$?\,\mathrm{slot~open}$}} & \shortstack[c]{0.6552852638634 \\ \scriptsize\textcolor{violet}{$?\,\mathrm{slot~open}$}} & \shortstack[c]{0.6552852639639 \\ \scriptsize\textcolor{violet}{$?\,\mathrm{slot~open}$}} \\
\parbox[t]{1.7in}{\raggedright DM half-slot candidate 2 (MeV)} & -- & \shortstack[c]{2.775832922367 \\ \scriptsize\textcolor{violet}{$?\,\mathrm{slot~open}$}} & \shortstack[c]{2.775832922379 \\ \scriptsize\textcolor{violet}{$?\,\mathrm{slot~open}$}} & \shortstack[c]{2.775832922805 \\ \scriptsize\textcolor{violet}{$?\,\mathrm{slot~open}$}} \\
\parbox[t]{1.7in}{\raggedright DM half-slot candidate 3 (MeV)} & -- & \shortstack[c]{7.267224937849 \\ \scriptsize\textcolor{violet}{$?\,\mathrm{slot~open}$}} & \shortstack[c]{7.267224937879 \\ \scriptsize\textcolor{violet}{$?\,\mathrm{slot~open}$}} & \shortstack[c]{7.267224938994 \\ \scriptsize\textcolor{violet}{$?\,\mathrm{slot~open}$}} \\
\parbox[t]{1.7in}{\raggedright DM half-slot candidate 4 (MeV)} & -- & \shortstack[c]{11.75861695333 \\ \scriptsize\textcolor{violet}{$?\,\mathrm{slot~open}$}} & \shortstack[c]{11.75861695338 \\ \scriptsize\textcolor{violet}{$?\,\mathrm{slot~open}$}} & \shortstack[c]{11.75861695518 \\ \scriptsize\textcolor{violet}{$?\,\mathrm{slot~open}$}} \\
\parbox[t]{1.7in}{\raggedright DM half-slot candidate 5 (MeV)} & -- & \shortstack[c]{30.78445884451 \\ \scriptsize\textcolor{violet}{$?\,\mathrm{slot~open}$}} & \shortstack[c]{30.78445884464 \\ \scriptsize\textcolor{violet}{$?\,\mathrm{slot~open}$}} & \shortstack[c]{30.78445884936 \\ \scriptsize\textcolor{violet}{$?\,\mathrm{slot~open}$}} \\
\parbox[t]{1.7in}{\raggedright DM half-slot candidate 6 (MeV)} & -- & \shortstack[c]{49.81030073569 \\ \scriptsize\textcolor{violet}{$?\,\mathrm{slot~open}$}} & \shortstack[c]{49.81030073590 \\ \scriptsize\textcolor{violet}{$?\,\mathrm{slot~open}$}} & \shortstack[c]{49.81030074354 \\ \scriptsize\textcolor{violet}{$?\,\mathrm{slot~open}$}} \\
\parbox[t]{1.7in}{\raggedright DM half-slot candidate 7 (MeV)} & -- & \shortstack[c]{130.4050603159 \\ \scriptsize\textcolor{violet}{$?\,\mathrm{slot~open}$}} & \shortstack[c]{130.4050603164 \\ \scriptsize\textcolor{violet}{$?\,\mathrm{slot~open}$}} & \shortstack[c]{130.4050603364 \\ \scriptsize\textcolor{violet}{$?\,\mathrm{slot~open}$}} \\
\parbox[t]{1.7in}{\raggedright DM half-slot candidate 8 (MeV)} & -- & \shortstack[c]{552.4047001080 \\ \scriptsize\textcolor{violet}{$?\,\mathrm{slot~open}$}} & \shortstack[c]{552.4047001104 \\ \scriptsize\textcolor{violet}{$?\,\mathrm{slot~open}$}} & \shortstack[c]{552.4047001951 \\ \scriptsize\textcolor{violet}{$?\,\mathrm{slot~open}$}} \\
\parbox[t]{1.7in}{\raggedright DM half-slot candidate 9 (MeV)} & -- & \shortstack[c]{0.6126241040672 \\ \scriptsize\textcolor{violet}{$?\,\mathrm{slot~open}$}} & \shortstack[c]{0.6126241040698 \\ \scriptsize\textcolor{violet}{$?\,\mathrm{slot~open}$}} & \shortstack[c]{0.6126241041638 \\ \scriptsize\textcolor{violet}{$?\,\mathrm{slot~open}$}} \\
\parbox[t]{1.7in}{\raggedright DM half-slot candidate 10 (MeV)} & -- & \shortstack[c]{1.732762633217 \\ \scriptsize\textcolor{violet}{$?\,\mathrm{slot~open}$}} & \shortstack[c]{1.732762633224 \\ \scriptsize\textcolor{violet}{$?\,\mathrm{slot~open}$}} & \shortstack[c]{1.732762633490 \\ \scriptsize\textcolor{violet}{$?\,\mathrm{slot~open}$}} \\
\parbox[t]{1.7in}{\raggedright DM half-slot candidate 11 (MeV)} & -- & \shortstack[c]{3.465525266434 \\ \scriptsize\textcolor{violet}{$?\,\mathrm{slot~open}$}} & \shortstack[c]{3.465525266449 \\ \scriptsize\textcolor{violet}{$?\,\mathrm{slot~open}$}} & \shortstack[c]{3.465525266980 \\ \scriptsize\textcolor{violet}{$?\,\mathrm{slot~open}$}} \\
\parbox[t]{1.7in}{\raggedright DM half-slot candidate 12 (MeV)} & -- & \shortstack[c]{9.801985665075 \\ \scriptsize\textcolor{violet}{$?\,\mathrm{slot~open}$}} & \shortstack[c]{9.801985665117 \\ \scriptsize\textcolor{violet}{$?\,\mathrm{slot~open}$}} & \shortstack[c]{9.801985666620 \\ \scriptsize\textcolor{violet}{$?\,\mathrm{slot~open}$}} \\
\parbox[t]{1.7in}{\raggedright DM half-slot candidate 13 (MeV)} & -- & \shortstack[c]{13.86210106574 \\ \scriptsize\textcolor{violet}{$?\,\mathrm{slot~open}$}} & \shortstack[c]{13.86210106580 \\ \scriptsize\textcolor{violet}{$?\,\mathrm{slot~open}$}} & \shortstack[c]{13.86210106792 \\ \scriptsize\textcolor{violet}{$?\,\mathrm{slot~open}$}} \\
\parbox[t]{1.7in}{\raggedright DM half-slot candidate 14 (MeV)} & -- & \shortstack[c]{27.72420213147 \\ \scriptsize\textcolor{violet}{$?\,\mathrm{slot~open}$}} & \shortstack[c]{27.72420213159 \\ \scriptsize\textcolor{violet}{$?\,\mathrm{slot~open}$}} & \shortstack[c]{27.72420213584 \\ \scriptsize\textcolor{violet}{$?\,\mathrm{slot~open}$}} \\
\multicolumn{5}{l}{\textit{RULED OUT BY FRAMEWORK}} \\
\parbox[t]{1.7in}{\raggedright 4th chiral generation} & \multicolumn{4}{l}{\parbox[t]{4.5in}{\raggedright\scriptsize\textcolor{red}{\textbf{ruled out:}} Framework predicts exactly three chiral fermion generations. \textbf{Falsified} if a 4th chiral lepton or quark generation is observed.}} \\
\parbox[t]{1.7in}{\raggedright Fundamental scalar beyond Higgs} & \multicolumn{4}{l}{\parbox[t]{4.5in}{\raggedright\scriptsize\textcolor{red}{\textbf{ruled out:}} Framework predicts the Higgs is the unique fundamental scalar. \textbf{Falsified} if a non-Higgs fundamental scalar is observed.}} \\
\parbox[t]{1.7in}{\raggedright Quantum-gravity scale below $\sim 10\,E_{\rm Planck}$} & \multicolumn{4}{l}{\parbox[t]{4.5in}{\raggedright\scriptsize\textcolor{red}{\textbf{ruled out:}} Framework predicts no quantum-gravity scale below $\sim 10\,E_{\rm Planck}$. \textbf{Falsified} if quantum-gravity signatures appear at sub-Planckian energies.}} \\
\parbox[t]{1.7in}{\raggedright Ultralight DM modes below $\sim 10^{-21}$~eV} & \multicolumn{4}{l}{\parbox[t]{4.5in}{\raggedright\scriptsize\textcolor{red}{\textbf{ruled out:}} Framework predicts no ultralight-DM modes below $\sim 10^{-21}$~eV (Lyman-$\alpha$ structure-formation cutoff, Rogers \& Peiris 2021). \textbf{Falsified} if direct detection finds a discrete-spectrum DM mode in this range.}} \\
\parbox[t]{1.7in}{\raggedright PQ axion (strong CP)} & \multicolumn{4}{l}{\parbox[t]{4.5in}{\raggedright\scriptsize\textcolor{red}{\textbf{ruled out:}} Framework predicts $\theta_{\rm QCD} = 0$ exactly with no PQ axion required. \textbf{Falsified} if $\theta_{\rm QCD}$ is measured nonzero.}} \\
\parbox[t]{1.7in}{\raggedright Dark photon (kinetic-mixing operator)} & \multicolumn{4}{l}{\parbox[t]{4.5in}{\raggedright\scriptsize\textcolor{red}{\textbf{ruled out:}} Framework predicts no 4D kinetic-mixing dark-photon operator. \textbf{Falsified} if dark-photon kinetic-mixing-style coupling is detected.}} \\
\end{longtable}
}

\end{widetext}

\section{Gravitational and cosmological observables}

\begin{widetext}
% Auto-generated by substrate_mass_spectrum.py (longtable)
{\scriptsize
\setlength{\LTcapwidth}{6.5in}
\begin{longtable}{lcccc}
\caption{Extended framework predictions: gravitational sector, cosmological observables, and additional constants. \textbf{Honest scope:} the framework's genuine gravity-sector predictions are exactly two --- Newton's $G$ ($\sim 2500$~ppb) and the Planck mass $m_P$ ($\sim 1$~ppb). All other gravity-sector entries (Schwarzschild radii, perihelion precessions, light deflection, Shapiro delay, gravitational redshift) are GR consequences computed using the framework's predicted $G$ together with \emph{observed} body parameters ($M_\odot$, $M_\oplus$, $R_\odot$, etc.); they are GR-consistency checks of the framework $G$, not independent predictions. The MOND $a_0$ and $\Lambda_{\rm crit-tension}$ rows are literature MOND (Milgrom 1983; Famaey \& Binney 2005) treated as ansatz; closure of the underlying equations is open. Computed at $\mathrm{mp.prec}=500$, anchored on three independent reading paths (CODATA $m_e$, CODATA $m_\mu$, framework Higgs vev). Each predicted value is shown to 15-digit precision with the deviation against the observed value beneath. The Observed column carries the published experimental uncertainty ($\pm$ ppm/ppb/ppt) below each value where available. $\dagger$ tension; $\ddagger$ open prediction; $\dagger\dagger$ regime-tension; $?$ pending closure (literature ansatz or GR-consistency check).}
\label{tab:extended-predictions} \\
\toprule
Observable & Observed & electron-anchor & muon-anchor & vev-anchor \\
\midrule
\endfirsthead
\multicolumn{5}{l}{\emph{Table~\ref{tab:extended-predictions} continued}} \\
\toprule
Observable & Observed & electron-anchor & muon-anchor & vev-anchor \\
\midrule
\endhead
\midrule
\multicolumn{5}{r}{\emph{(continued on next page)}} \\
\endfoot
\bottomrule
\endlastfoot
\parbox[t]{1.7in}{\raggedright Higgs vev (structural identity) (GeV)} & \shortstack[c]{246.2196500000 \\ \scriptsize$\pm 24.368~\mathrm{ppm}$} & \shortstack[c]{246.2196500197 \\ \scriptsize\textcolor{OliveGreen}{$\checkmark\,+80.025~\mathrm{ppt}$}} & \shortstack[c]{246.2196500197 \\ \scriptsize\textcolor{OliveGreen}{$\checkmark\,+80.025~\mathrm{ppt}$}} & \shortstack[c]{246.2196500197 \\ \scriptsize\textcolor{OliveGreen}{$\checkmark\,+80.025~\mathrm{ppt}$}} \\
\parbox[t]{1.7in}{\raggedright Newton's $G$ (m$^3$kg$^{-1}$s$^{-2}$)} & \shortstack[c]{6.674300000000e-11 \\ \scriptsize$\pm 22.474~\mathrm{ppm}$} & \shortstack[c]{6.674316918458e-11 \\ \scriptsize\textcolor{OliveGreen}{$\checkmark\,+2.535~\mathrm{ppm}$}} & \shortstack[c]{6.674316918401e-11 \\ \scriptsize\textcolor{OliveGreen}{$\checkmark\,+2.535~\mathrm{ppm}$}} & \shortstack[c]{6.674316916354e-11 \\ \scriptsize\textcolor{OliveGreen}{$\checkmark\,+2.535~\mathrm{ppm}$}} \\
\parbox[t]{1.7in}{\raggedright Planck mass $m_P$ (GeV)} & \shortstack[c]{1.220890000000e+19 \\ \scriptsize$\pm 11.467~\mathrm{ppm}$} & \shortstack[c]{1.220888581931e+19 \\ \scriptsize\textcolor{OliveGreen}{$\checkmark\,-1.162~\mathrm{ppm}$}} & \shortstack[c]{1.220888581936e+19 \\ \scriptsize\textcolor{OliveGreen}{$\checkmark\,-1.161~\mathrm{ppm}$}} & \shortstack[c]{1.220888582124e+19 \\ \scriptsize\textcolor{OliveGreen}{$\checkmark\,-1.161~\mathrm{ppm}$}} \\
\parbox[t]{1.7in}{\raggedright GR perihelion precession (Mercury) ($\,\mathrm{arcsec/century}$)} & \shortstack[c]{43.11000000000 \\ \scriptsize$\pm 4871.260~\mathrm{ppm}$} & \shortstack[c]{42.99184861659 \\ \scriptsize\textcolor{OliveGreen}{$\checkmark\,-2740.696~\mathrm{ppm}$}} & \shortstack[c]{42.99184861659 \\ \scriptsize\textcolor{OliveGreen}{$\checkmark\,-2740.696~\mathrm{ppm}$}} & \shortstack[c]{42.99184861659 \\ \scriptsize\textcolor{OliveGreen}{$\checkmark\,-2740.696~\mathrm{ppm}$}} \\
\parbox[t]{1.7in}{\raggedright GR perihelion precession (Venus) ($\,\mathrm{arcsec/century}$)} & \shortstack[c]{8.626000000000 \\ \scriptsize$\pm 3477.858~\mathrm{ppm}$} & \shortstack[c]{8.626741339097 \\ \scriptsize\textcolor{OliveGreen}{$\checkmark\,+85.942~\mathrm{ppm}$}} & \shortstack[c]{8.626741339097 \\ \scriptsize\textcolor{OliveGreen}{$\checkmark\,+85.942~\mathrm{ppm}$}} & \shortstack[c]{8.626741339097 \\ \scriptsize\textcolor{OliveGreen}{$\checkmark\,+85.942~\mathrm{ppm}$}} \\
\parbox[t]{1.7in}{\raggedright GR perihelion precession (Earth) ($\,\mathrm{arcsec/century}$)} & \shortstack[c]{3.840000000000 \\ \scriptsize$\pm 13020.833~\mathrm{ppm}$} & \shortstack[c]{3.839649170448 \\ \scriptsize\textcolor{OliveGreen}{$\checkmark\,-91.362~\mathrm{ppm}$}} & \shortstack[c]{3.839649170448 \\ \scriptsize\textcolor{OliveGreen}{$\checkmark\,-91.362~\mathrm{ppm}$}} & \shortstack[c]{3.839649170448 \\ \scriptsize\textcolor{OliveGreen}{$\checkmark\,-91.362~\mathrm{ppm}$}} \\
\parbox[t]{1.7in}{\raggedright GR perihelion precession (Mars) ($\,\mathrm{arcsec/century}$)} & \shortstack[c]{1.350000000000 \\ \scriptsize$\pm 37037.037~\mathrm{ppm}$} & \shortstack[c]{1.351262711937 \\ \scriptsize\textcolor{OliveGreen}{$\checkmark\,+935.342~\mathrm{ppm}$}} & \shortstack[c]{1.351262711937 \\ \scriptsize\textcolor{OliveGreen}{$\checkmark\,+935.342~\mathrm{ppm}$}} & \shortstack[c]{1.351262711937 \\ \scriptsize\textcolor{OliveGreen}{$\checkmark\,+935.342~\mathrm{ppm}$}} \\
\parbox[t]{1.7in}{\raggedright GR perihelion precession (Jupiter) ($\,\mathrm{arcsec/century}$)} & \shortstack[c]{0.06230000000000 \\ \scriptsize$\pm 16051.364~\mathrm{ppm}$} & \shortstack[c]{0.06232950659433 \\ \scriptsize\textcolor{OliveGreen}{$\checkmark\,+473.621~\mathrm{ppm}$}} & \shortstack[c]{0.06232950659433 \\ \scriptsize\textcolor{OliveGreen}{$\checkmark\,+473.621~\mathrm{ppm}$}} & \shortstack[c]{0.06232950659433 \\ \scriptsize\textcolor{OliveGreen}{$\checkmark\,+473.621~\mathrm{ppm}$}} \\
\parbox[t]{1.7in}{\raggedright Light deflection (solar limb) ($\,\mathrm{arcsec}$)} & \shortstack[c]{1.751700000000 \\ \scriptsize$\pm 456.699~\mathrm{ppm}$} & \shortstack[c]{1.751644036032 \\ \scriptsize\textcolor{OliveGreen}{$\checkmark\,-31.948~\mathrm{ppm}$}} & \shortstack[c]{1.751644036032 \\ \scriptsize\textcolor{OliveGreen}{$\checkmark\,-31.948~\mathrm{ppm}$}} & \shortstack[c]{1.751644036032 \\ \scriptsize\textcolor{OliveGreen}{$\checkmark\,-31.948~\mathrm{ppm}$}} \\
\parbox[t]{1.7in}{\raggedright Shapiro round-trip (Earth-Saturn conj., $b = 1.6\,R_\odot$) ($\,\mu\mathrm{s}$)} & \shortstack[c]{260.0000000000 \\ \scriptsize$\pm 76923.077~\mathrm{ppm}$} & \shortstack[c]{264.4693472015 \\ \scriptsize\textcolor{OliveGreen}{$\checkmark\,+17189.797~\mathrm{ppm}$}} & \shortstack[c]{264.4693472015 \\ \scriptsize\textcolor{OliveGreen}{$\checkmark\,+17189.797~\mathrm{ppm}$}} & \shortstack[c]{264.4693472015 \\ \scriptsize\textcolor{OliveGreen}{$\checkmark\,+17189.797~\mathrm{ppm}$}} \\
\parbox[t]{1.7in}{\raggedright Solar grav. redshift $z$ (limb)} & \shortstack[c]{2.120000000000e-6 \\ \scriptsize$\pm 4716.981~\mathrm{ppm}$} & \shortstack[c]{2.123052482755e-6 \\ \scriptsize\textcolor{OliveGreen}{$\checkmark\,+1439.850~\mathrm{ppm}$}} & \shortstack[c]{2.123052482755e-6 \\ \scriptsize\textcolor{OliveGreen}{$\checkmark\,+1439.850~\mathrm{ppm}$}} & \shortstack[c]{2.123052482755e-6 \\ \scriptsize\textcolor{OliveGreen}{$\checkmark\,+1439.850~\mathrm{ppm}$}} \\
\parbox[t]{1.7in}{\raggedright Schwarzschild radius (Sun) ($\,\mathrm{km}$)} & \shortstack[c]{2.953000000000 \\ \scriptsize$\pm 338.639~\mathrm{ppm}$} & \shortstack[c]{2.954015224506 \\ \scriptsize\textcolor{orange}{$\bullet\,+343.794~\mathrm{ppm}$}} & \shortstack[c]{2.954015224506 \\ \scriptsize\textcolor{orange}{$\bullet\,+343.794~\mathrm{ppm}$}} & \shortstack[c]{2.954015224506 \\ \scriptsize\textcolor{orange}{$\bullet\,+343.794~\mathrm{ppm}$}} \\
\parbox[t]{1.7in}{\raggedright Schwarzschild radius (Earth) ($\,\mathrm{mm}$)} & \shortstack[c]{8.870000000000 \\ \scriptsize$\pm 112.740~\mathrm{ppm}$} & \shortstack[c]{8.870125356371 \\ \scriptsize\textcolor{OliveGreen}{$\checkmark\,+14.133~\mathrm{ppm}$}} & \shortstack[c]{8.870125356371 \\ \scriptsize\textcolor{OliveGreen}{$\checkmark\,+14.133~\mathrm{ppm}$}} & \shortstack[c]{8.870125356371 \\ \scriptsize\textcolor{OliveGreen}{$\checkmark\,+14.133~\mathrm{ppm}$}} \\
\parbox[t]{1.7in}{\raggedright MOND scale $a_0 = cH_0/(2\pi)$ ($\,\mathrm{m/s^2}$)} & \shortstack[c]{$1.2\!\times\!10^{-10}$ \\ \scriptsize$\pm 0.02\;\mathrm{stat}$ \\ \scriptsize$\pm 0.24\;\mathrm{syst}$} & \shortstack[c]{$1.0422\!\times\!10^{-10}$ (CMB) \\ $1.1288\!\times\!10^{-10}$ (local) \\ \scriptsize\textcolor{blue}{$?$\,Milgrom 1983;}\\ \scriptsize\textcolor{blue}{ansatz, open}} & \shortstack[c]{$1.0422\!\times\!10^{-10}$ (CMB) \\ $1.1288\!\times\!10^{-10}$ (local) \\ \scriptsize\textcolor{blue}{$?$\,Milgrom 1983;}\\ \scriptsize\textcolor{blue}{ansatz, open}} & \shortstack[c]{$1.0422\!\times\!10^{-10}$ (CMB) \\ $1.1288\!\times\!10^{-10}$ (local) \\ \scriptsize\textcolor{blue}{$?$\,Milgrom 1983;}\\ \scriptsize\textcolor{blue}{ansatz, open}} \\
\parbox[t]{1.7in}{\raggedright $\Lambda_{\rm crit-tension} = (H_0/c)^2$ ($\,\mathrm{m^{-2}}$)} & 1.097121526139e-52 & \shortstack[c]{$5.3086e-53$ (CMB) \\ $6.2273e-53$ (local) \\ \scriptsize\textcolor{blue}{$?$\,Friedmann factor}\\ \scriptsize\textcolor{blue}{$3\Omega_\Lambda{\sim}2$ open}} & \shortstack[c]{$5.3086e-53$ (CMB) \\ $6.2273e-53$ (local) \\ \scriptsize\textcolor{blue}{$?$\,Friedmann factor}\\ \scriptsize\textcolor{blue}{$3\Omega_\Lambda{\sim}2$ open}} & \shortstack[c]{$5.3086e-53$ (CMB) \\ $6.2273e-53$ (local) \\ \scriptsize\textcolor{blue}{$?$\,Friedmann factor}\\ \scriptsize\textcolor{blue}{$3\Omega_\Lambda{\sim}2$ open}} \\
\parbox[t]{1.7in}{\raggedright Hubble constant $H_0$ (km/s/Mpc)} & \shortstack[c]{$\sim 73$ (local) \\ \scriptsize$\sim 67$ (CMB)} & \shortstack[c]{\textcolor{blue}{open} \\ \scriptsize\textcolor{blue}{$\dagger\dagger$ regime} \\ \scriptsize\textcolor{blue}{tension; pending}} & \shortstack[c]{\textcolor{blue}{open} \\ \scriptsize\textcolor{blue}{$\dagger\dagger$ regime} \\ \scriptsize\textcolor{blue}{tension; pending}} & \shortstack[c]{\textcolor{blue}{open} \\ \scriptsize\textcolor{blue}{$\dagger\dagger$ regime} \\ \scriptsize\textcolor{blue}{tension; pending}} \\
\end{longtable}
}

\end{widetext}

\section{Predicted but unobserved (falsification targets)}

\noindent\textbf{Discrete dark-matter spectrum (21 modes total: 13
with specific mass predictions, 8 high-mass with envelope).}
The framework predicts a discrete spectrum of dark-matter masses,
none yet detected. The 13 specific-mass predictions are:
\begin{itemize}\itemsep0pt
\item \emph{Ultralight band (7 modes):} 0.62~$\mu$eV, 2.23~neV,
8.03~peV, 28.9~feV, 0.10~aeV, 0.37~zeV, 1.34~zeV. The
Lyman-$\alpha$-forest ultralight-DM bound $m \gtrsim 2\times
10^{-20}$~eV at 95\%~CL~\cite{RogersPeiris2021} excludes the
1.34~zeV mode; the next-lightest mode (0.37~zeV is below; 0.10~aeV
$\sim 19\times$ above the bound) is above the bound.
\item \emph{MeV--GeV band (6 modes):} 0.83~MeV, 5.83~MeV,
46.6~MeV, 1.49~GeV, 4.82~GeV, 593~GeV.
\end{itemize}
The remaining 8 high-mass modes lie within the envelope
$[191~\mathrm{GeV},\,1.1~\mathrm{EeV}]$; specific masses
$\ddagger$~open.

\noindent\textbf{Slot-occupancy disclaimer.}
The dark-matter modes enumerated above are \emph{candidate slots}
only. The framework predicts the mass that would be carried by
each slot \emph{if} it is occupied; it does not yet predict which
slots are physically populated, nor with what relic abundance. We
are explicitly \emph{not} predicting that detector experiments will
find a particle in every listed mass band. Some, all, or none of
these candidate slots may host observable dark-matter species; the
framework's slot-occupancy prediction function (which slots are
populated and at what abundance) is unfinished work and is gated
by the open J-items below. The 14 spin-$\tfrac12$ half-integer-slot
modes (auto-flagged violet in the predictions table) are the
clearest example of this: each has a definite predicted mass, but
whether the half-integer slot can support a populated state at all
is a spin-cover question the framework cannot currently answer.
Treat the dark-matter spectrum as a list of \emph{permissible}
masses for any dark-sector state the framework does admit, not as
a list of guaranteed detections.

\noindent\textbf{CP rationals (framework predictions).}
$\delta_{\rm CP}^{\rm PMNS} = 11\pi/10 \approx 198^\circ$ (consistent
with NuFit-6.0~\cite{NuFit60} normal ordering: at the IC24 fit
$\delta_{\rm CP}^{\rm best} = 212^\circ$ with asymmetric range
$-41^\circ/+26^\circ$, the framework's $198^\circ$ sits at
$\sim 0.3\sigma$ on the low-$\delta$ side; awaiting DUNE/HyperK).
$\delta_{\rm CKM} = 11\pi/30 \approx 1.152$~rad (within $\sim 1\sigma$
of the CKMfitter/PDG~2024 central value~\cite{CKMfitter2005,PDG2024}).

\noindent\textbf{Three matter generations ($N_{\rm gen} = 3$).}
The framework predicts exactly three chiral fermion generations.

\noindent\textbf{Strong CP: $\theta_{\rm QCD} = 0$ exactly}, with no
PQ axion required.

\noindent\textbf{Internal topology-discrimination test: $|t| = 6.92\sigma$.}
A measurement-level discrimination test against the leading
alternative scenario clears $5\sigma$ on both percentile and BCa
$95\%$-CI floors.

\noindent\textbf{No quantum-gravity scale below $\sim 10\,E_{\rm Planck}$;
no fourth chiral generation; no fundamental scalar beyond Higgs.}

\section{Upcoming experiments and what we predict they will find}

The following named experiments have data releases or sensitivity
upgrades scheduled in the 2026--2030 window that directly test the
framework's predictions. Each entry pairs the experiment with the
specific framework prediction it probes and the central-value
expectation we anticipate the data to converge on.

\subsection*{Neutrino sector}

\paragraph{JUNO (Jiangmen Underground Neutrino Observatory).}
\emph{Status:} taking data since August 2025. \emph{First-result
window:} 2026--2030. The 59-day exposure already published
$\sin^2\theta_{12} = 0.3092 \pm 0.0087$ and
$\Delta m_{21}^2 = (7.50\pm 0.12)\times 10^{-5}~\mathrm{eV}^2$;
the full 6-year program will reach $\sim\!0.5\%$ precision on
$\sin^2\theta_{12}$. \emph{Framework prediction:}
$\sin^2\theta_{12} = 0.3066$. The current measurement sits at
$\sim\!0.3\sigma$ above the framework value; we expect JUNO's
high-statistics result to converge on a central value within
$\sim\!0.5\%$ of $0.3066$.

\paragraph{Hyper-Kamiokande.}
\emph{Status:} under construction; commissioning 2027--2028; data
taking from 2028. \emph{First-result window:} 2029--2030.
\emph{Framework prediction:} $\delta_{\rm CP}^{\rm PMNS} = 11\pi/10
\approx 198^\circ$ (normal ordering). Hyper-K's projected $6^\circ$--$20^\circ$
precision on $\delta_{\rm CP}$ will discriminate $198^\circ$ from
the current global-fit best fits ($177^\circ$ in NuFit-6.0~IC19,
$212^\circ$ in IC24) at high significance. We anticipate Hyper-K
will converge near $198^\circ$.

\paragraph{IceCube-Upgrade.}
\emph{Status:} 7 new strings deploying austral summer 2025--2026.
\emph{First-result window:} 2027--2028. \emph{Framework prediction:}
$\sin^2\theta_{23} = 4/7 \approx 0.5714$. We anticipate
IceCube-Upgrade's atmospheric-oscillation result to land near
$0.5714$.

\paragraph{T2K + NOvA joint analysis.}
\emph{Status:} first joint result published October 2025; next
combined update expected 2027--2028. The current $3\sigma$ band on
$\delta_{\rm CP}$ does not exclude CP conservation. \emph{Framework
prediction:} same $11\pi/10$ as Hyper-K. We anticipate the next
T2K$+$NOvA joint update will pull the central value toward $198^\circ$
as data accumulates.

\subsection*{CKM sector}

\paragraph{LHCb (Run 3 $\to$ Upgrade II).}
\emph{Status:} Run~3 data ongoing; analysis 2026--2028. The
current measurement $\gamma = (64.6 \pm 2.8)^\circ$ is already
within $0.5\sigma$ of the framework's $11\pi/30 \approx 66.0^\circ$.
\emph{First-result window:} Run 3 datasets give $\sim\!1^\circ$
precision by 2027--2028; HL-LHC$+$Upgrade~II projects $0.3^\circ$
by $\sim\!2036$. We anticipate LHCb's central value to converge
near $66^\circ$.

\paragraph{Belle II.}
\emph{Status:} running. Independent $\gamma$ extraction expected
$\sim\!2027$ providing cross-check on LHCb. \emph{Framework
prediction:} same $66^\circ$. Independent confirmation expected.

\subsection*{Cosmology and CMB}

\paragraph{Simons Observatory.}
\emph{Status:} Large Aperture Telescope first light February 2025;
two Small Aperture Telescopes running since 2023; full array
online mid-2026. \emph{First-result window:} 2027--2028.
\emph{Framework prediction:} $n_s = 58/60 \approx 0.9667$. SO's
projected $\sigma(n_s) \approx 0.002$ (vs Planck's $0.0042$) will
discriminate $0.9667$ from the current Planck-2018 central value
$0.9649 \pm 0.0042$ at greater than $2\sigma$. We anticipate SO's
central value to land near $0.9667$.

\paragraph{LiteBIRD.}
\emph{Status:} JAXA satellite; launch JFY 2027 ($\sim\!2028$);
3-year survey. Joint constraint on $n_s$ and tensor-to-scalar
ratio $r$ at $\sigma(r) \sim 10^{-3}$. Confirms or refutes the
$n_s$ prediction in conjunction with SO.

\paragraph{Vera C. Rubin Observatory (LSST).}
\emph{Status:} operations begin early 2026; alert stream live
February 2026; DR1 $\sim\!2028$. SN-Ia$+$lensing program will
deliver an independent $H_0$ at $\sim\!1\%$ precision by 2028--2030.
\emph{Framework prediction:} the local-vs-CMB $H_0$ tension
($\sim\!9\%$) is a genuine substrate-physical effect, not
measurement error; Rubin's local $H_0$ should land near the
SH0ES central value, $\sim\!73~\mathrm{km/s/Mpc}$.

\paragraph{Euclid \& DESI.}
\emph{Status:} both running; Euclid DR1 expected 2026--2027; DESI
5-year results expected 2027 with extension through 2028.
\emph{Framework prediction:} BAO/lensing-derived $H_0$ should
remain in the CMB-extrapolated band ($\sim\!67$--$68~\mathrm{km/s/Mpc}$);
neither survey will resolve the $\sim\!9\%$ tension to within
measurement error. The persistence of the gap is the framework
prediction.

\paragraph{CMB-S4.}
\emph{Status:} revised plan published June 2025; full survey starts
mid-2030s. Long-term ultimate test of $n_s = 0.9667$ at $\sigma(r)
\sim 5\times 10^{-4}$ tensor sensitivity.

\subsection*{Dark-matter direct detection}

\paragraph{LZ (LUX-ZEPLIN).}
\emph{Status:} operating; latest world-best WIMP exclusion limits
published December 2025. Next major release $\sim\!2027$.
\emph{Framework prediction:} discrete-mass dark-matter peaks at
$1.49~\mathrm{GeV}$, $4.82~\mathrm{GeV}$, and $593~\mathrm{GeV}$ are
within LZ's accessible window. We expect LZ to either detect a
peak at one of these specific masses (consistent with the
framework) or continue to set null limits across the band.

\paragraph{XENONnT.}
\emph{Status:} running since 2021; updates 2026--2028. Same
framework predictions apply; XENONnT provides an independent
cross-check.

\paragraph{XLZD (DARWIN successor).}
\emph{Status:} Design Book published 2025; construction late 2020s;
operations early 2030s. 60--80~tonne LXe target reaches
$\sigma_{\rm SI} \sim 3\times 10^{-49}~\mathrm{cm}^2$, covering the
framework's MeV--GeV and high-mass-envelope predictions down to
the neutrino fog.

\paragraph{ADMX (axion direct detection).}
\emph{Status:} currently scanning $1.10$--$1.31~\mathrm{GHz}$
($4.6$--$5.4~\mu\mathrm{eV}$); upgrades extend the range through
2030. \emph{Framework prediction:} no QCD axion (since $\theta_{\rm QCD}
= 0$ is a structural identity, not a Peccei-Quinn relaxation), so
we predict ADMX will continue to set null limits and
\emph{not} detect any DFSZ-class axion signal in any frequency
band.

\paragraph{NANOGrav, EPTA, SKA (pulsar-timing arrays).}
\emph{Status:} NANOGrav 15-year dataset published; SKA-Mid first
light 2028; 30-year datasets simulated through the 2030s.
\emph{Framework prediction:} the framework's $0.37~\mathrm{zeV}$ and
$1.34~\mathrm{zeV}$ ultralight modes lie in the PTA-sensitive
$10^{-23}$--$10^{-22}~\mathrm{eV}$ range. We anticipate
PTA-program detection of one of these specific zeV-band masses if
they are dominant DM components, or continued null limits.

\subsection*{Lepton magnetic moments}

\paragraph{J-PARC E34 muon $g{-}2$/EDM.}
\emph{Status:} beam commissioning 2025; data taking $\sim\!2030$;
target precision $450~\mathrm{ppb}$. Independent ultracold-muon
technique with different systematics. Cross-checks the closed
Fermilab Muon $g{-}2$ result; framework predicts no discrepancy
beyond Standard-Model lattice-QCD evaluations.

\paragraph{Belle II tau program.}
\emph{Status:} running. Few-${\times}10^{-6}$ sensitivity on $a_\tau$
projected by 2028.

\subsection*{Higgs and electroweak precision}

\paragraph{HL-LHC (ATLAS \& CMS).}
\emph{Status:} Run 3 finishes $\sim\!2026$; HL-LHC Run 4 begins
mid-2030s; 3000~$\mathrm{fb}^{-1}$ per detector. Reaches
$\kappa_b \sim 3\%$, $\kappa_\tau \sim 2\%$, total $\Gamma_H \sim 5\%$
precision on Higgs partial widths. \emph{Framework predictions:}
the predicted partial widths and the observation that no fourth
chiral fermion generation exists below $\sim\!2.5~\mathrm{TeV}$;
we expect HL-LHC to confirm the Standard-Model partial-width
pattern and to push 4th-generation exclusion to $\sim\!2.5~\mathrm{TeV}$
without discovery. Di-Higgs $\to \lambda_{HHH}$ extraction at
$3$--$5\sigma$ tests the no-new-fundamental-scalar prediction.

\subsection*{Gravitational tests}

\paragraph{LIGO/Virgo/KAGRA O5 \& LISA.}
\emph{Status:} O5 run starting 2027; LISA launch $\sim\!2035$.
Multi-messenger BNS$+$GRB timing constrains $\Delta c/c$ down to
$10^{-17}$ and tests vacuum dispersion of gravitational waves.
\emph{Framework prediction:} no frame-dependent $c$, no vacuum
dispersion; we anticipate continued null measurement of any
deviation from the speed-of-light universality.

\paragraph{ACES (Atomic Clock Ensemble in Space).}
\emph{Status:} launched April 21, 2025 to ISS. Tests gravitational
redshift to $\sim\!3\times 10^{-6}$ and time variation of fundamental
constants. \emph{Framework prediction:} no time variation of
constants beyond the framework's structural integers, no
Lorentz-violation signature.

\paragraph{LHC searches for sub-Planckian quantum gravity.}
\emph{Status:} ongoing through HL-LHC. \emph{Framework prediction:}
no detection of quantum-gravity signatures (e.g.\ micro-black-hole
production) below $\sim\!10\,E_{\rm Planck}$.

\subsection*{Anchored summary of near-term decisive tests}

The most actionable falsification targets within the 2026--2028
window are:
\begin{itemize}\itemsep0pt
\item \textbf{Simons Observatory} on $n_s = 0.9667$ (2027--2028
result, $>2\sigma$ discrimination from Planck central);
\item \textbf{JUNO} on $\sin^2\theta_{12} = 0.3066$ (already at
$0.3\sigma$ in 59-day data; full-precision result by $\sim\!2030$);
\item \textbf{LHCb Run 3} on $\gamma = 11\pi/30$ (already at
$0.5\sigma$; $\sim\!1^\circ$ precision by 2027--2028);
\item \textbf{Hyper-Kamiokande} on $\delta_{\rm CP}^{\rm PMNS} =
11\pi/10$ (first decisive results 2029--2030);
\item \textbf{LZ/XENONnT} on the MeV--GeV dark-matter spectrum
(continuous updates through 2027).
\end{itemize}

\section{What would falsify the framework}

The list below is mixed: items marked $\dagger$ are committed and
falsifiable today; items marked $\ddagger$ become falsifiable once
upstream prediction work closes (see \S~Work remaining).

\begin{enumerate}\itemsep0pt
\item[$\dagger$] Fourth chiral fermion generation observed.
\item[$\dagger$] $\theta_{\rm QCD}$ measured nonzero.
\item[$\ddagger$] Direct-detection dark-matter peak at a mass
either (a) outside the union of the 13 specific-mass predictions
$\pm$ predicted width AND outside the high-mass envelope
$[191~\mathrm{GeV},\,1.1~\mathrm{EeV}]$, or (b) inside the high-mass
envelope but inconsistent with the 8 high-mass values once those
are committed.
\item[$\ddagger$] $\delta_{\rm CP}^{\rm PMNS}$ measured inconsistent
with $11\pi/10$ at $\geq 3\sigma$ in normal ordering. (Currently
auto-flagged blue; becomes a committed-falsifiable prediction once
the upstream PMNS-mixing work closes.)
\item[$\dagger$] Topology-discrimination null result: a reproducible
measurement-level test showing the framework's preferred scenario
loses to the leading alternative at $> 5\sigma$.
\item[$\dagger$] Astrophysical detection of an azimuthal mode-pattern
on a rotating compact body (e.g.\ pulsar surface oscillations via
NICER or successor X-ray timing missions) whose symmetry group,
classified as a finite subgroup of $O(3)$, falls outside
$\{\mathbb{Z}_n, D_n, T, O, I\}$.
\item[$\ddagger$] Stable non-collapsed compact object whose mass
exceeds the framework-predicted substrate TOV bound $M_{\rm sub-TOV}$
(value $\ddagger$~open; bound is geometric at saturation,
$r_{\rm phys} = r_s$).
\item[$\dagger$] CMB-S4 measurement of $n_s$ inconsistent with the
framework's prediction $n_s = (N-2)/N$ at the $\Delta n_s\sim 10^{-4}$
level, with $N=60$ as the framework canonical value
($n_s = 58/60 = 0.9667$ at $N=60$ vs.\ $n_s = 56/58 = 0.9655$ at
$N=58$ differ by $\sim 1.2\times 10^{-3}$, so the $N$-dependence
is itself part of the falsifiable claim).
\item[$\dagger$] Frame-dependent $c$, vacuum dispersion, or any
departure from speed-of-light universality at any energy scale.
\item[$\dagger$] Experimental confirmation of a quantum-gravity
scale below $\sim 10\,E_{\rm Planck} \approx 10^{20}$~GeV (via
gravitational decoherence at observer-energies $\ll E_{\rm Planck}$,
or any single-graviton-exchange process detected at sub-Planckian
energies).
\end{enumerate}

\section{Work remaining (prioritized by unblock count)}

The two tables together carry 38 distinct blue-flagged
observables that are \emph{not yet falsifiable} because the
upstream prediction work that would commit them at structural
precision is not yet closed. The list below tallies open work
items in priority order, counting how many blue observables each
item turns into committed predictions.

\paragraph{J1. Lepton/quark mixing structure ---
unblocks 12 observables.}
Blocks: PMNS $\sin^2\theta_{12}$, $\sin^2\theta_{23}$,
$\sin^2\theta_{13}$, $\delta_{\rm CP}$ (4); CKM
$\sin^2\theta_{12}^{\rm CKM}$, $\sin^2\theta_{23}^{\rm CKM}$,
$\sin^2\theta_{13}^{\rm CKM}$, $\delta_{\rm CKM}$ (4); Higgs partial
widths $\Gamma_{H\to bb,\tau\tau,cc,ff}$ (4).

\paragraph{J2. Gauge-coupling precision --- unblocks 9 observables.}
Blocks: $\alpha_{\rm EM}^{\rm tree}$, $\alpha_2$, $\alpha_1^{\rm EW}$,
$\alpha_1^{\rm GUT}$ (4); $g_2$, $g_1^{\rm EW}$, $g_1^{\rm GUT}$ (3);
$\sin^2\theta_W$ (1); $Z$ boson cascading from the
$W$-mass$+\sin^2\theta_W$ identity (1).

\paragraph{J3. $\alpha_{\rm EM}$ hadronic running ---
unblocks 5 observables.} Blocks: lepton anomalous magnetic
moments $a_e^{\rm LO}$, $a_\mu^{\rm LO}$, $a_\tau^{\rm LO}$ (3);
Bohr radius $a_0$ (1); Rydberg energy $\mathrm{Ry}$ (1).

\paragraph{J4. Gravity / cosmology block --- unblocks 4
observables.} A single closure decomposes into five sub-items:
(a) the form of the MOND interpolation function (cf.\
Famaey--Binney~\cite{FamaeyBinney2005} and the radial-acceleration
relation~\cite{McGaughLelliSchombert2016}); (b) the prefactor in
the relation $a_0 = c H_0/(2\pi)$ (cf.\ Milgrom~1983
\cite{Milgrom1983}); (c) the residual $\sim 50\%$ Friedmann
coefficient that closes $\Lambda_{\rm crit-tension} = (H_0/c)^2$
to the full $\Lambda$CDM relation; (d) the local-vs-CMB $H_0$
regime tension at $\sim 9\%$ (SH0ES~\cite{Riess2022}
vs.\ Planck~\cite{Planck2018VI}); (e) the dark-mode spatial
density at cluster scales needed for combined-MOND-plus-dark-mode
acceleration tests. Blocks: $\rho_\Lambda$ (dark-tower contribution);
MOND $a_0$; $\Lambda_{\rm crit-tension}$; $H_0$ regime tension.

\paragraph{J5. NLO electroweak radiative corrections ---
unblocks 4 observables.} Blocks: $\Gamma_W$, $\Gamma_Z$,
$\Gamma_Z^{\rm invisible}$, $\Gamma_Z^{\rm leptonic}$ (4).

\paragraph{J6. NLO$+$ QED/QCD corrections for lepton lifetimes
--- unblocks 2 observables.} Blocks: $\tau_\mu$, $\tau_\tau$ (2).

\paragraph{J7. $\Lambda_{\rm QCD}$ closure --- unblocks 1
observable.} Blocks: $\alpha_s$ at lower scales (1).

\paragraph{J8. Thermal-history closure --- unblocks 1 observable.}
Blocks: $N_{\rm eff}$ at recombination (1).

\paragraph{Total unblock count.}
J1 (12) $+$ J2 (9) $+$ J3 (5) $+$ J4 (4) $+$ J5 (4) $+$ J6 (2)
$+$ J7 (1) $+$ J8 (1) $=$ 38 observables, matching the count of
distinct blue-flagged observables across both tables.

\paragraph{Strategic note.}
J1 and J2 together account for 21 of 38 blue observables (55\%);
J3 and J7 are tightly coupled. Realistically, closing J1$+$J2$+$J3$+$J7
together unblocks 27 of 38 (71\%); J4 is structurally independent
and is the second-priority closure.

\paragraph{$\ddagger$-flagged items.} $\rho_\Lambda$ and the
$H_0$ regime tension fall under J4 (sub-items~c and~d). The TOV
bound's nuclear-scale closure is a separate item not yet in the
J1--J8 tally; it remains $\ddagger$-flagged.

\section{Provenance}

All numerical content auto-generated at $\mathrm{mp.prec}=500$.
Source code, per-section reproducibility traces, and JSON output
artefacts mirrored at
\href{https://github.com/gable-digital/papers}{github.com/gable-digital/papers}.

\begin{thebibliography}{99}

\bibitem{CKMfitter2005}
J.~Charles \emph{et al.} (CKMfitter Group),
``CP violation and the CKM matrix: assessing the impact of the asymmetric
$B$ factories,''
\emph{Eur.~Phys.~J.~C} \textbf{41}, 1--131 (2005),
\href{https://doi.org/10.1140/epjc/s2005-02169-1}{doi:10.1140/epjc/s2005-02169-1},
\href{https://arxiv.org/abs/hep-ph/0406184}{arXiv:hep-ph/0406184}.

\bibitem{NuFit60}
I.~Esteban, M.~C.~Gonzalez-Garcia, M.~Maltoni, I.~Martinez-Soler,
J.~P.~Pinheiro, and T.~Schwetz,
``NuFit-6.0: updated global analysis of three-flavor neutrino oscillations,''
\emph{J.~High Energy Phys.} \textbf{2024}, 216 (2024),
\href{https://doi.org/10.1007/JHEP12(2024)216}{doi:10.1007/JHEP12(2024)216},
\href{https://arxiv.org/abs/2410.05380}{arXiv:2410.05380}.

\bibitem{CODATA2022}
P.~J.~Mohr, D.~B.~Newell, B.~N.~Taylor, and E.~Tiesinga,
``CODATA recommended values of the fundamental physical constants: 2022,''
\emph{Rev.~Mod.~Phys.} \textbf{97}, 025002 (2025),
\href{https://doi.org/10.1103/RevModPhys.97.025002}{doi:10.1103/RevModPhys.97.025002},
\href{https://arxiv.org/abs/2409.03787}{arXiv:2409.03787}.

\bibitem{FamaeyBinney2005}
B.~Famaey and J.~Binney,
``Modified Newtonian dynamics in the Milky Way,''
\emph{Mon.~Not.~R.~Astron.~Soc.} \textbf{363}, 603--608 (2005),
\href{https://doi.org/10.1111/j.1365-2966.2005.09474.x}{doi:10.1111/j.1365-2966.2005.09474.x},
\href{https://arxiv.org/abs/astro-ph/0506723}{arXiv:astro-ph/0506723}.

\bibitem{McGaughLelliSchombert2016}
S.~S.~McGaugh, F.~Lelli, and J.~M.~Schombert,
``The radial acceleration relation in rotationally supported galaxies,''
\emph{Phys.~Rev.~Lett.} \textbf{117}, 201101 (2016),
\href{https://doi.org/10.1103/PhysRevLett.117.201101}{doi:10.1103/PhysRevLett.117.201101},
\href{https://arxiv.org/abs/1609.05917}{arXiv:1609.05917}.

\bibitem{Milgrom1983}
M.~Milgrom,
``A modification of the Newtonian dynamics as a possible alternative to
the hidden mass hypothesis,''
\emph{Astrophys.~J.} \textbf{270}, 365--370 (1983),
\href{https://doi.org/10.1086/161130}{doi:10.1086/161130}.

\bibitem{PDG2024}
S.~Navas \emph{et al.} (Particle Data Group),
``Review of particle physics,''
\emph{Phys.~Rev.~D} \textbf{110}, 030001 (2024),
\href{https://doi.org/10.1103/PhysRevD.110.030001}{doi:10.1103/PhysRevD.110.030001}.

\bibitem{Planck2018VI}
N.~Aghanim \emph{et al.} (Planck Collaboration),
``Planck 2018 results.~VI.~Cosmological parameters,''
\emph{Astron.~Astrophys.} \textbf{641}, A6 (2020),
\href{https://doi.org/10.1051/0004-6361/201833910}{doi:10.1051/0004-6361/201833910},
\href{https://arxiv.org/abs/1807.06209}{arXiv:1807.06209}.

\bibitem{Riess2022}
A.~G.~Riess \emph{et al.},
``A comprehensive measurement of the local value of the Hubble constant
with $1~\mathrm{km\,s^{-1}\,Mpc^{-1}}$ uncertainty from the Hubble Space
Telescope and the SH0ES Team,''
\emph{Astrophys.~J.~Lett.} \textbf{934}, L7 (2022),
\href{https://doi.org/10.3847/2041-8213/ac5c5b}{doi:10.3847/2041-8213/ac5c5b},
\href{https://arxiv.org/abs/2112.04510}{arXiv:2112.04510}.

\bibitem{RogersPeiris2021}
K.~K.~Rogers and H.~V.~Peiris,
``Strong bound on canonical ultra-light axion dark matter from the
Lyman-alpha forest,''
\emph{Phys.~Rev.~Lett.} \textbf{126}, 071302 (2021),
\href{https://doi.org/10.1103/PhysRevLett.126.071302}{doi:10.1103/PhysRevLett.126.071302},
\href{https://arxiv.org/abs/2007.12705}{arXiv:2007.12705}.

\end{thebibliography}

\end{document}
